\documentclass[12pt,a4paper]{article}
\usepackage[T1]{fontenc}
\usepackage[utf8]{inputenc}
\usepackage[ngerman]{babel}
\usepackage{amsmath,amssymb,amsthm}
\usepackage{geometry}
\geometry{margin=2.5cm}
\usepackage{mathrsfs}

\newcommand{\tang}{\tan}

\newtheorem*{satz}{Satz}
\newtheorem*{theorem}{Theorem}

\title{Lehrbuch der Algebra --- Band 1, Teil 14}
\author{Heinrich Weber}
\date{}

\begin{document}
\maketitle

\noindent\textbf{\S. 173. Primideale in $\Omega_m$}\\[4pt]
\noindent\textit{(Fortsetzung)}

\bigskip\noindent und wenn wir dies auf jeden Factor von $F_n(t)$ anwenden, und beachten, dass $p^{h+1}$ nach dem Modul $m$ dieselbe Zahlenreihe durchläuft, wie $p^h$, so folgt:
\[
[F_n(t)]^p \equiv F_n(t^p) \pmod{p},
\]
und diese Congruenz gilt dann natürlich auch für den Modul $p$.

Dies ist aber das Kennzeichen dafür, dass $F_n(t)$ mit einem Polynom $P_n(t)$ mit ganzen rationalen Zahlencoëfficienten nach dem Modul $p$ congruent ist (\S. 150, 4.), also
\begin{equation}\label{eq:4}
F_n(t) \equiv P_n(t) \pmod{p}.
\end{equation}

Setzen wir dies in (3) ein, so ergiebt sich zunächst eine Congruenz nach dem Modul $p$, die aber, da es eine Congruenz zwischen rationalen Functionen ist, auch für den Modul $p$ bestehen muss, also
\begin{equation}\label{eq:5}
f(t) \equiv P_{\xi_1}(t)\, P_{\xi_2}(t) \ldots P_{\xi_e}(t) \pmod{p}.
\end{equation}

Nach der Definition (1), (\ref{eq:4}) ist
\begin{equation}\label{eq:6}
P_1(r) \equiv 0 \pmod{p},
\end{equation}
und die sämmtlichen Wurzeln dieser Congruenz sind
\begin{equation}\label{eq:7}
r,\; r^p,\; \ldots,\; r^{p^{f-1}}.
\end{equation}

Ebenso hat die Congruenz
\[
P_n(t) \equiv 0 \pmod{p}
\]
die Wurzeln
\[
r^n,\; r^{np},\; \ldots,\; r^{np^{f-1}},
\]
und diese Wurzeln sind, wenn man $n$ das System der Zahlen $\xi_1, \xi_2, \ldots, \xi_e$ durchlaufen lässt, alle incongruent nach dem Modul $p$.

Macht man in der Congruenz
\[
F_1(t) \equiv P_1(t) \pmod{p}
\]
die Substitution $(r, r^n)$, so geht $F_1(t)$ in $F_n(t)$, $p$ in das conjugirte Ideal $p_n$ über, und $P_1(t)$ bleibt als rationale Form ungeändert. Wir haben also
\begin{equation}\label{eq:8}
F_n(t) \equiv P_1(t) \pmod{p_n},
\end{equation}
und wenn wir hierin $t = r$ setzen, so folgt:
\begin{equation}\label{eq:9}
F_n(r) \equiv P_1(r) \pmod{p_n},
\end{equation}
$F_n(r)$ ist aber, wenn $n$ nicht in $\mathfrak{A}$ enthalten ist, relativ prim zu $p$, also nicht durch $p_n$ theilbar.

Es ist also $P_1(r)$ nach (9) durch $p_1$, aber durch keinen der mit $p_1$ conjugirten Primfactoren theilbar, und es folgt, da $p$ nicht durch $p^2$ theilbar ist:

\begin{satz}
1. Der Primfactor $p_1$ ist der grösste gemeinschaftliche Theiler von $p$ und $P_1(r)$; in gleicher Weise ergiebt sich, dass $p_n$ der grösste gemeinschaftliche Theiler von $p$ und $P_1(r^n)$ ist, oder auch der grösste gemeinschaftliche Theiler von $p$ und $P_{n'}(r)$, wenn $nn' \equiv 1 \pmod{m}$ ist.
\end{satz}

Das letztere ergiebt sich aus der aus (\ref{eq:4}) durch die Substitution $(r, r^{n'})$ folgenden Congruenz
\[
F_1(t) \equiv P_n(t) \pmod{p_{n'}}.
\]

Man erhält also die sämmtlichen Primfactoren von $p$, wenn man die grössten gemeinschaftlichen Theiler von $p$ mit jeder der rationalen Functionen
\begin{equation}\label{eq:10}
P_{\xi_1}(r),\; P_{\xi_2}(r),\; \ldots,\; P_{\xi_e}(r)
\end{equation}
aufsucht, und zwar ist $p_{\xi_i}$ der grösste gemeinschaftliche Theiler von $p$ und $P_{\xi_i}^{-1}$.

Wir wollen den Fall noch etwas näher betrachten, wo $f = 1$ ist, also
\begin{equation}\label{eq:11}
p \equiv 1 \pmod{m}.
\end{equation}

In diesem Falle ist $e = \varphi(m)$; die Gruppe $\mathfrak{A}$ reducirt sich auf die Einheit, und die sämmtlichen $\mu$ conjugirten Factoren $p$ von $p$ sind von einander verschieden. Die Formen
\[
P_{\xi_1}(t),\; P_{\xi_2}(t),\; \ldots,\; P_{\xi_e}(t)
\]
werden alle linear, d.\,h.\ $r$ ist einer rationalen Zahl nach jedem der Moduln $p_i$ congruent.

Dies ergiebt sich auch daraus, dass hier die Anzahl der nach dem Modul $p$ incongruenten Zahlen gleich $N(p) = p$ ist, und dass also $0, 1, 2, \ldots, p-1$ ein volles Restsystem ist.

Ist hiernach etwa $r \equiv c \pmod{p}$, so muss, da $r^m = 1$ ist, $c^m \equiv 1 \pmod{p}$ also auch $(\bmod\, p)$ sein, und wenn also $g$ eine primitive Wurzel der Primzahl $p$ ist, so muss
\begin{equation}\label{eq:12}
c \equiv g^{-n\frac{p-1}{m}} \pmod{p}
\end{equation}
sein, worin $n$ relativ prim zu $m$ ist. Lassen wir $p$ das System der conjugirten Ideale durchlaufen, so muss $n$ in (12) die Gruppe $\mathfrak{N}$ durchlaufen, und es ist also nur Sache der Bezeichnung, wenn wir festsetzen:
\begin{equation}\label{eq:13}
r \equiv g^{-\frac{p-1}{m}} \pmod{p}.
\end{equation}

\bigskip\noindent\textbf{\S. 174. Das Kummer'sche Theorem}\\[4pt]

Machen wir darin die Substitution $(r, r^n)$, so wird
\begin{equation}\label{eq:14}
r^n \equiv g^{-n\frac{p-1}{m}} \pmod{p_n},
\end{equation}
oder wenn wir $n'$ durch die Congruenz
\begin{equation}\label{eq:15}
nn' \equiv 1 \pmod{m}
\end{equation}
definiren:
\begin{equation}\label{eq:16}
r \equiv g^{-n'\frac{p-1}{m}} \pmod{p_n}.
\end{equation}

Es ist also in diesem Falle, übereinstimmend mit der allgemeinen Regel, $p_n$ der grösste gemeinschaftliche Theiler von $p$ und $\bigl(r - g^{-n'\frac{p-1}{m}}\bigr)$.

Durch diese Bestimmung sind die einzelnen Primformen $p_n$ genau charakterisirt. Wir erhalten also den Satz:

\begin{satz}
2. Eine Primzahl $p$, die nach dem Modul $m$ mit $1$ congruent ist, zerfällt im Körper $\Omega_m$ in $\varphi(m)$ von einander verschiedene Primfactoren $p_n$ vom ersten Grade, die man als grösste gemeinschaftliche Theiler von $p$ mit den verschiedenen Zahlen $r - g^{-n'\frac{p-1}{m}}$ erhält, wenn $g$ eine primitive Wurzel von $p$ ist und $n'$ durch die Congruenz $nn' \equiv 1 \pmod{m}$ bestimmt ist.
\end{satz}

Nach \S. 141 können wir jede ganze oder gebrochene Zahl $\omega$ und jedes Functional des Körpers $\Omega_m$ in der Weise in Primfactoren zerlegen, dass Zähler und Nenner keinen gemeinschaftlichen Primfactor enthalten, und diese Zerlegung ist, von Einheitsfactoren abgesehen, völlig bestimmt, so dass, wenn $p$ ein Primideal ist, ein ganz bestimmter positiver oder negativer Exponent $k$ existirt, so dass das Functional $\omega\, p^{-k}$ den Primfactor $k$ weder im Zähler noch im Nenner enthält. Wir sagen dann der Kürze halber, es sei $p^k$ die in $\omega$ enthaltene Potenz von $p$.

Ist $k = 0$, so sagen wir, $p$ ist in $\omega$ nicht enthalten.

Zwei Zahlen (oder auch Functionale) $\omega, \omega'$ eines Körpers, in denen ein und dasselbe Primfunctional $p$ enthalten ist, wollen wir theilerverwandt oder kurz verwandt nennen; wenn dagegen kein Primfunctional zugleich in $\omega$ und in $\omega'$ enthalten ist, so heissen $\omega$ und $\omega'$ theilerfremd oder fremd. Besonders nützlich wird uns dieser Ausdruck in der Folge sein, wenn an Stelle von $\omega'$ die natürliche Primzahl $p$ tritt, die durch $p$ theilbar ist, und wir reden danach von den mit einer Zahl $\omega$ verwandten oder zu ihr fremden Primzahlen.

Eine Zahl $\omega$, die gar keine verwandten Primzahlen hat, ist eine Einheit.

Im Körper $\Omega_m$ (wie überhaupt in jedem Normalkörper) sind alle conjugirten Zahlen $\omega_n$ mit denselben Primzahlen verwandt.

\bigskip\noindent\textbf{\S. 175. Die Einheitswurzeln im Körper $\Omega_m$}\\[4pt]

Für die weiter zu machenden Anwendungen ist eine genauere Kenntniss der Einheiten des Körpers $\Omega_m$ erforderlich, die ja auch an sich von grossem Interesse ist. Wir beschäftigen uns hier nicht mit den functionalen, sondern nur mit den numerischen Einheiten, worunter, wie wir uns erinnern, ganze Zahlen des Körpers $\Omega_m$ zu verstehen sind, deren Norm $\pm 1$ ist, oder eine ganze Zahl, deren reciproker Werth gleichfalls ganz ist.

Zu den Einheiten gehören sicher alle in $\Omega_m$ enthaltenen Einheitswurzeln; wir können aber leicht beweisen, dass diese Einheitswurzeln nur Potenzen von $r$ sein können, mit positivem und negativem Vorzeichen.

Es sei nämlich $\varrho$ eine in $\Omega_m$ enthaltene Einheitswurzel, und
\begin{equation}\label{eq:175-1}
\varrho = \varphi(r).
\end{equation}
Ist $q^h$ die höchste im Grade von $\varrho$ aufgehende Potenz von $q$, so ist $\varrho^{q^h}$ eine Einheitswurzel, deren Grad nicht durch $q$ theilbar ist, und
\begin{equation}\label{eq:175-2}
\varrho^{q^h} = [\varphi(r)]^{q^h}.
\end{equation}

Wegen der Irreducibilität der Kreistheilungsgleichung im weiteren Sinne\footnote{Vgl. Bd. I, \S. 134 und den Nachtrag zu diesem Bande.} können daher in (2) die sämmtlichen Substitutionen $(r, r^n)$ gemacht werden, ohne dass die linke Seite sich ändert, und folglich ist $\varrho^{q^h}$ eine rationale Zahl, und da es eine Einheit ist, muss es $= \pm 1$ sein. Mithin ist $\varrho$, vom Vorzeichen abgesehen, einer Einheitswurzel vom Grade $q^h$ oder, wenn $q = 2$ ist, vom Grade $q^{h+1}$ gleich. Es ist nachzuweisen, dass $q^h$ oder $q^{h+1}$ nicht grösser als $m$ sein kann. Dies ergiebt sich aber aus (1). Denn wäre $\varrho$ eine Einheitswurzel höheren als $m^{\text{ten}}$ Grades, so wäre $r$ eine Potenz von $\varrho$, und $\varrho$ müsste einer irreduciblen Gleichung höheren Grades als $r$, also auch höheren Grades als $\varphi(r)$ genügen, was nach (1) ein Widerspruch ist. Also haben wir den Satz:

\begin{satz}
1. Die einzigen in $\Omega_m$ enthaltenen Einheitswurzeln sind die mit positivem und negativem Zeichen genommenen Potenzen von $r$.
\end{satz}

Hieran schliesst sich ein ganz allgemeiner, d.\,h.\ in jedem algebraischen Zahlkörper gültiger Satz über die Einheitswurzeln, den wir jetzt beweisen wollen.

Der Satz lautet:

\begin{satz}
2. Ist $\alpha$ eine ganze Zahl eines algebraischen Zahlkörpers $n^{\text{ten}}$ Grades, und haben alle mit $\alpha$ conjugirten Zahlen $\alpha_1, \alpha_2, \ldots, \alpha_n$ den absoluten Werth $1$, so ist $\alpha$ eine Einheitswurzel\footnote{Kronecker, „Zwei Sätze über Gleichungen mit ganzzahligen Coëfficienten''. Crelle's Journal, Bd. 53 (1857). Minkowski, „Geometrie der Zahlen'', Art. 43.}.
\end{satz}

Zum Beweis ist zunächst zu bemerken, dass der absolute Werth der Norm einer ganzen Zahl $\omega$ dem Product der absoluten Werthe der conjugirten Zahlen $\omega_1, \omega_2, \ldots, \omega_n$ gleich, und als ganze rationale Zahl jedenfalls grösser oder gleich $1$ ist. Es ist daher nicht möglich, dass alle diese absoluten Werthe kleiner als $1$ sind. Sind sie aber alle gleich $1$, wie bei der in unserem Satze angenommenen Zahl $\alpha$, so muss die Norm $= \pm 1$ sein, und $\alpha$ ist eine Einheit. Ist $\beta$ eine zweite Einheit, die mit ihren conjugirten zugleich den absoluten Werth $1$ hat, so hat der Quotient $\alpha : \beta$ dieselbe Eigenschaft. Wenn daher unter den zu diesem Quotienten conjugirten Zahlen auch nur eine reell ist, so muss
\[
\frac{\alpha}{\beta} = \pm 1
\]
oder $\alpha = \pm\beta$ sein, und dies gilt dann auch noch, wenn $\alpha, \beta$ durch die entsprechenden Zahlen eines conjugirten Körpers ersetzt werden.

Wenn also $\alpha$ und $\beta$ weder gleich noch entgegengesetzt sind, so ist ihr Verhältniss nicht reell, und es besteht daher zwischen den absoluten Werthen die Ungleichung:
\[
|\alpha \pm \beta| < |\alpha| + |\beta|.
\]
(Vergl. die Einleitung zum ersten Bande, S. 19.)

Es ist also
\[
|\alpha \pm \beta| < 2,
\]
und folglich kann $\frac{1}{2}(\alpha \pm \beta)$ keine ganze Zahl sein, weil der absolute Werth der Norm dieser Zahl kleiner als $1$ ist. Wenn nun $\omega_1, \omega_2, \ldots, \omega_n$ eine Basis von $\mathfrak{o}$ ist, und
\begin{align*}
\alpha &= a_1\omega_1 + a_2\omega_2 + \cdots + a_n\omega_n\\
\beta &= b_1\omega_1 + b_2\omega_2 + \cdots + b_n\omega_n,
\end{align*}
so können die ganzen rationalen Zahlen $a, b$ nicht den Congruenzen
\[
a_1 \equiv b_1,\; a_2 \equiv b_2,\; \ldots,\; a_n \equiv b_n \pmod{2}
\]
genügen, weil sonst $\frac{1}{2}(\alpha - \beta)$ eine ganze Zahl wäre.

Wenn wir also die sämmtlichen Zahlen $\alpha$ in $2^n$ Fächer vertheilen, indem wir alle Zahlen in ein Fach werfen, in denen $a_1, a_2, \ldots, a_n$ dieselben Reste (0 oder 1) nach dem Modul 2 lassen, so können in jedem dieser Fächer höchstens 2 Zahlen, nämlich $\alpha$ und $-\alpha$, vorkommen, und wenn wir $\alpha = 0$ noch ausschliessen, so giebt es sogar in einem dieser Fächer, in dem die $a_1, a_2, \ldots, a_n$ alle gerade sind, gar keine Zahl $\alpha$. Die Anzahl aller möglichen Zahlen $\alpha$ ist also endlich und höchstens $= 2^{n+1} - 2$.

Wenn nun der absolute Werth von $\alpha = 1$ ist, so gilt dasselbe von allen Potenzen von $\alpha$. Folglich muss in der unbegrenzten Reihe
\[
1,\; \alpha,\; \alpha^2,\; \alpha^3,\; \ldots
\]
nothwendig dieselbe Zahl zum zweiten Male wiederkehren, also $\alpha^h = \alpha^k$ und $k > h$ sein. Dann ist aber
\[
\alpha^{k-h} = 1,
\]
d.\,h.\ $\alpha$ ist eine Einheitswurzel, wie bewiesen werden sollte.

Wir fügen noch die Bemerkung bei, dass, wenn eine Zahl $\alpha$ eines Körpers $\Omega$ eine Einheitswurzel $m^{\text{ten}}$ Grades ist, auch alle mit $\alpha$ conjugirten Zahlen Einheitswurzeln desselben Grades sind. Denn in der rationalen Gleichung $\alpha^m - 1 = 0$ kann $\alpha$ durch jeden der conjugirten Werthe ersetzt werden.

Die Anzahl der Einheitswurzeln, die in einem Körper $\Omega$ enthalten sind, kann immer nur eine endliche sein, weil jede Zahl in $\Omega$ einer rationalen Gleichung genügt, deren Grad dem Körpergrad höchstens gleich ist. Der Grad der Einheitswurzeln in $\Omega$ kann daher einen endlichen Werth nicht übersteigen.

\bigskip\noindent\textbf{\S. 176. Der in $\Omega_m$ enthaltene reelle Körper $H_m$}\\[4pt]

Jede Zahl des Körpers $\Omega_m$ geht durch die Substitution $(r, r^{-1})$ in die conjugirt imaginäre Zahl über, die auch durch die Vertauschung $(i, -i)$ erhalten wird. Das Product zweier solcher conjugirt imaginärer Zahlen ist das Quadrat des absoluten Werthes einer jeden von ihnen.

Eine Zahl in $\Omega_m$ ist reell, wenn sie durch die Vertauschung $(r, r^{-1})$ ungeändert bleibt, und nur unter dieser Voraussetzung. Jede solche Zahl lässt sich rational durch die zweigliedrige Periode $r + r^{-1}$ ausdrücken, und gehört also einem reellen Körper vom Grade $\frac{1}{2}\varphi(m)$ an, der ein Theiler von $\Omega_m$ ist. Diesen wollen wir den in $\Omega_m$ enthaltenen reellen Körper nennen und mit $H_m$ bezeichnen (obwohl auch noch andere reelle Körper, nämlich alle Theiler von $H_m$, in $\Omega_m$ enthalten sind).

Die $\frac{1}{2}\varphi(m) = \frac{1}{2}\mu$ Zahlen
\begin{equation}\label{eq:176-1}
1,\; r + r^{-1},\; r^2 + r^{-2},\; \ldots,\; r^{\frac{1}{2}\mu-1} + r^{-\frac{1}{2}\mu+1}
\end{equation}
bilden eine Basis der ganzen Zahlen des Körpers $H_m$. Denn nach \S. 170, (7) ist
\begin{equation}\label{eq:176-2}
\begin{aligned}
\omega = x_0 &+ x_1 r + x_2 r^2 + \cdots + x_{\frac{1}{2}\mu-1} r^{\frac{1}{2}\mu-1} + x_{\frac{1}{2}\mu} r^{\frac{1}{2}\mu}\\
&+ x_{-1} r^{-1} + x_{-2} r^{-2} + \cdots + x_{-\frac{1}{2}\mu+1} r^{-\frac{1}{2}\mu+1}
\end{aligned}
\end{equation}
dann und nur dann eine ganze Zahl des Körpers $\Omega_m$, wenn die Coëfficienten $x_0, x_1, \ldots$ ganze rationale Zahlen sind.

Die Bedingung dafür, dass diese Zahl dem Körper $H_m$ angehört, erhält man, wenn man die Substitution $(r, r^{-1})$ ausführt und die Differenz $= 0$ setzt. Dividirt man noch durch $r - r^{-1}$, so ergiebt sich so:
\begin{align*}
&(x_1 - x_{-1}) + (x_2 - x_{-2})(r + r^{-1}) + \cdots\\
&\quad + (x_{\frac{1}{2}\mu-1} - x_{-\frac{1}{2}\mu+1}) \frac{r^{\frac{1}{2}\mu-1} - r^{-\frac{1}{2}\mu+1}}{r - r^{-1}}\\
&\quad + x_{\frac{1}{2}\mu} \frac{r^{\frac{1}{2}\mu} - r^{-\frac{1}{2}\mu}}{r - r^{-1}} = 0.
\end{align*}

Die Divisionen lassen sich hier ausführen, und wenn man dann mit $r^{\frac{1}{2}\mu-1}$ multiplicirt, so ergiebt sich für $r$ eine rationale Gleichung von niedrigerem als $\mu^{\text{ten}}$ Grade, die nur dann befriedigt sein kann, wenn alle ihre Coëfficienten verschwinden. Dies führt zu den Gleichungen
\[
x_{\frac{1}{2}\mu} = 0,\; x_{\frac{1}{2}\mu-1} = x_{-\frac{1}{2}\mu+1},\; \ldots,\; x_1 = x_{-1},
\]
und es ist also jede ganze Zahl $\omega$ des Körpers $H_m$ in der Form enthalten:
\begin{align*}
\omega = x_0 &+ x_1(r + r^{-1}) + x_2(r^2 + r^{-2}) + \cdots\\
&+ x_{\frac{1}{2}\mu-1}(r^{\frac{1}{2}\mu-1} + r^{-\frac{1}{2}\mu+1}).
\end{align*}

Statt der Basis (1) kann man auch, wenn man
\[
\varrho = r + r^{-1}
\]
setzt, die Potenzen von $\varrho$:
\begin{equation}\label{eq:176-3}
1,\; \varrho,\; \varrho^2,\; \ldots,\; \varrho^{\frac{1}{2}\mu-1}
\end{equation}
als Basis der ganzen Zahlen von $H_m$ wählen. Denn die Grössen (3) lassen sich ganzzahlig durch die (1) ausdrücken und umgekehrt.

Die Zahl $\varrho$ ist die Wurzel einer irreduciblen Gleichung vom Grade $\frac{1}{2}\mu$, die wir mit $\psi(\varrho) = 0$ bezeichnen wollen. Ist $f(r) = 0$ die Gleichung für $r$ (\S. 169), so besteht die identische Relation
\begin{equation}\label{eq:176-4}
f(x) = x^{\frac{1}{2}\mu}\,\psi\Bigl(x + \frac{1}{x}\Bigr),
\end{equation}
woraus durch Bildung der Ableitung
\begin{equation}\label{eq:176-5}
f'(r) = r^{\frac{1}{2}\mu}\,\psi'(\varrho)\,(1 - r^{-2}).
\end{equation}

Hieraus können wir die Grundzahl $\Delta_1$ des Körpers $H_m$ bilden, die, da $H_m$ nur reelle Zahlen enthält, positiv sein muss. Diese Grundzahl ist, wenn wir die Norm in Bezug auf $H_m$ mit $N_1$ bezeichnen:
\[
\Delta_1 = \pm N_1\psi'(\varrho).
\]

Ist aber $N$ die in Bezug auf $\Omega_m$ gebildete Norm, so ist
\[
N\psi'(\varrho) = [N_1\psi'(\varrho)]^2 = \Delta_1^2,
\]
und demnach ergiebt die Formel (5) mit Rücksicht auf \S. 169:
\[
\Delta = N(1 - r^{-2})\,\Delta_1^2.
\]

Nach \S. 169, (7) und (10) ist:
\begin{alignat*}{3}
N(1 - r^{-2}) &= q &\quad &\text{bei ungeradem } q\\
&= 4 &\quad &\text{für } q = 2,\\
\pm\Delta &= q\,\Delta_1^2 &\quad &\text{bei ungeradem } q\\
&= 4\,\Delta_1^2 &\quad &\text{für } q = 2.
\end{alignat*}

Setzt man darin
\[
\Delta = \pm\, q^{q^{\varkappa-1}[\varkappa(q-1)-1]}
\]
so folgt
\begin{equation}\label{eq:176-7}
\begin{aligned}
\Delta_1 &= q^{q^{\varkappa}-1\frac{q-1}{2} - \frac{q^{\varkappa-1}+1}{2}} \quad \text{bei ungeradem } q\\
&= 2^{(\varkappa-1)2^{\varkappa}-2-1} \qquad\qquad\; \text{für } q = 2,
\end{aligned}
\end{equation}
so dass $\Delta_1$ immer eine Potenz von $q$ ist.

\bigskip\noindent\textbf{\S. 177. Die Primideale im Körper $H_m$}\\[4pt]

Es ist jetzt die Frage zu untersuchen, in welcher Beziehung die Primideale des Körpers $H_m$ zu denen des Körpers $\Omega_m$ stehen.

Wenn $\mathfrak{P}$ irgend ein Primideal in $H_m$ vom Grade $f_1$ bedeutet, so muss $\mathfrak{P}$ Theiler einer natürlichen Primzahl $p$ sein, und $\mathfrak{P}$ kann also (nach \S. 171) nur dann durch die zweite oder eine höhere Potenz eines Primfactors $p$ in $\Omega_m$ theilbar sein, wenn $p = q$ ist.

Ist aber $\mathfrak{P}$ ein Theiler von $q$, so muss es eine Potenz von $\sigma = 1 - r$ sein. Wir setzen etwa
\[
\mathfrak{P} = \sigma^\lambda.
\]

Nehmen wir hiervon die Norm in Bezug auf $\Omega_m$, die das Quadrat der Norm in Bezug auf $H_m$ ist, so folgt (nach \S. 169)
\[
q^{2f_1} = q^\lambda,
\]
also $f_1 = \frac{1}{2}\lambda$. Da $f_1$ eine ganze Zahl ist, so muss $\lambda$ mindestens $= 2$ sein. Es kann aber $\lambda$ auch nicht grösser als $2$ sein, da $\sigma^2$ mit der in $H_m$ enthaltenen ganzen Zahl $(1-r)(1-r^{-1})$ associirt ist, also $\sigma^2$ durch $\mathfrak{P}$ theilbar sein muss. Demnach ist $f_1 = 1$, und wir erhalten den Satz:

\begin{satz}
1. Die Primzahl $q$ ist die $\frac{1}{2}\mu^{\text{te}}$ Potenz einer in $H_m$ existirenden Primzahl ersten Grades.
\end{satz}

Es sei nun $p$ von $q$ verschieden, und $p$ ein Primfactor von $\mathfrak{P}$ im Körper $\Omega_m$ vom Grade $f$, der durch die Substitution $(r, r^{-1})$ in $p'$ übergeht. Dann ist $\mathfrak{P}$ sowohl durch $p$ als durch $p'$ theilbar, und andererseits ist $pp'$ in $H_m$ enthalten und also durch $\mathfrak{P}$ theilbar, und wir haben zu unterscheiden, ob $p, p'$ verschieden sind oder nicht.

Ist $p$ von $p'$ verschieden, so ist $\mathfrak{P}$ durch $pp'$ theilbar, und es folgt $\mathfrak{P} = pp'$, und durch Normbildung $f_1 = f$.

Ist aber $p = p'$, so ist $\mathfrak{P} = p$, und die Normbildung ergiebt $2f_1 = f$.

Welcher dieser beiden Fälle eintritt, das hängt also davon ab, ob das Primideal $p$ die Substitution $(r, r^{-1})$ gestattet oder nicht, d.\,h.\ ob $(r, r^{-1})$ in der Gruppe, zu der das Primideal gehört, enthalten ist oder nicht. Nach \S. 172 ist dieser Unterschied dadurch bedingt, ob es irgend einen Exponenten $h$ giebt, für den die Congruenz
\[
p^h \equiv -1 \pmod{m}
\]
erfüllt ist, oder ob es keinen solchen Exponenten giebt.

Wir fassen das Ergebniss folgendermaassen zusammen:

\begin{satz}
2. Die von $q$ verschiedenen natürlichen Primzahlen $p$ zerfallen in zwei Arten $p_1, p_2$.

Die erste Art $p_1$ ist dadurch characterisirt, dass für irgend einen Exponenten $h$
\[
p_1^h \equiv -1 \pmod{m},
\]
und diese Primzahlen zerfallen, wenn sie für den Modul $m$ zum Exponenten $f$ gehören und $\mu = ef$ gesetzt wird, im Körper $H_m$ in $e$ Primfactoren $\frac{1}{2}f^{\text{ten}}$ Grades. Ihre Primfactoren sind auch in $\Omega_m$ nicht weiter zerlegbar.

Die Primzahlen der zweiten Art $p_2$ sind dadurch bestimmt, dass keine ihrer Potenzen nach dem Modul $m$ mit $-1$ congruent wird, und sie zerfallen im Körper $H_m$ in $\frac{1}{2}e$ Primfactoren $f^{\text{ten}}$ Grades. Ihre Primfactoren sind in $\Omega_m$ noch in je zwei Factoren zerlegbar.
\end{satz}

Bei den Primzahlen der ersten Art muss $f$ sicher eine gerade Zahl sein. Wenn $m$ ungerade ist, so genügt es auch, damit $p$ eine Primzahl von der ersten Art sei, dass sie zu einem geraden Exponenten gehöre; denn dann ist
\[
(p^{\frac{1}{2}f} - 1)(p^{\frac{1}{2}f} + 1) \equiv 0 \pmod{m}.
\]

Der erste dieser Factoren $p^{\frac{1}{2}f} - 1$ ist aber nicht durch $m$ theilbar, weil sonst $p$ zum Exponenten $\frac{1}{2}f$ gehören würde. Folglich muss $p^{\frac{1}{2}f} + 1$ durch $m$ theilbar sein. Hier können aber nicht beide Factoren $p^{\frac{1}{2}f} - 1$ und $p^{\frac{1}{2}f} + 1$ durch $q$ theilbar sein, weil sonst auch ihre Differenz, die $= 2$ ist, durch $q$ theilbar wäre, und folglich ist $p^{\frac{1}{2}f} + 1$ durch $m$ theilbar.

Ist aber $m$ eine Potenz von $2$, so ist die Congruenz
\[
p^h \equiv -1 \pmod{m}
\]
nur dann möglich, wenn $p \equiv -1 \pmod{m}$ ist. Denn jede gerade Potenz einer ungeraden Zahl ist $\equiv +1 \pmod{8}$ und kann also nicht $\equiv -1 \pmod{m}$ sein. Ist aber $h$ ungerade, so ist
\[
\frac{p^h + 1}{p + 1} = p^{h-1} - p^{h-2} + \cdots + 1
\]
selbst ungerade, und folglich ist $p^h + 1$ durch keine höhere Potenz von $2$ theilbar, als $p + 1$. Demnach können wir dem Satze 2. noch folgende nähere Bestimmung hinzufügen:

\begin{satz}
3. Wenn $m$ ungerade ist, so gehören die Primzahlen $p_1$ zu einem geraden, die Primzahlen $p_2$ zu einem ungeraden Exponenten.

Ist $m$ gerade, so gehören die Primzahlen der Form $km - 1$ zur ersten und alle anderen ungeraden Primzahlen zur zweiten Art.
\end{satz}

\bigskip\noindent\textbf{\S. 178. Die Einheiten des Körpers $H_m$}\\[4pt]

Die dem Körper $H_m$ angehörigen Einheiten sind von besonderer Wichtigkeit für die Anwendungen. Auf sie lassen sich auch die Einheiten des Körpers $\Omega_m$ zurückführen nach folgendem Satze:

\begin{satz}
1. Jede Einheit des Körpers $\Omega_m$ ist das Product einer Einheitswurzel und einer reellen Einheit des Körpers $\Omega_m$.
\end{satz}

Bezeichnen wir irgend eine Einheit des Körpers $\Omega_m$ mit $\mathcal{E}(r)$, so ist $\mathcal{E}(r^{-1})$ der conjugirt imaginäre Werth, der gleichfalls eine Einheit darstellt, und der Quotient $\mathcal{E}(r) : \mathcal{E}(r^{-1})$ ist wieder eine Einheit, deren absoluter Werth
\begin{equation}\label{eq:178-1}
\sqrt{\frac{\mathcal{E}(r)}{\mathcal{E}(r^{-1})}\,\frac{\mathcal{E}(r^{-1})}{\mathcal{E}(r)}}
\end{equation}
gleich $1$ ist, und folglich ist dieser Quotient eine Einheitswurzel und mithin $= \pm r^h$, worin $h$ irgend ein ganzzahliger Exponent ist (\S. 175, Satz 1., 2.). Wir haben also
\begin{equation}\label{eq:178-2}
\mathcal{E}(r) = \pm r^h\, \mathcal{E}(r^{-1}).
\end{equation}

Es ist nun im weiteren Beweise ein kleiner Unterschied zu machen, je nachdem $m$ ungerade oder gerade ist.

Ist $m$ zunächst ungerade, so können wir in (2) den Exponenten $h$ gerade annehmen, da wir ihn nöthigenfalls durch $h + m$ ersetzen können. Ferner ist in der Formel (2) nur das obere Zeichen zulässig. Denn wenn das untere Zeichen stände, so würde folgen:
\[
\mathcal{E}(r) \equiv \mathcal{E}(1) \equiv -\mathcal{E}(1),\quad 2\,\mathcal{E}(1) \equiv 0 \pmod{(1-r)}.
\]

Nach \S. 169 ist aber $1 - r$ ein Theiler von $q$, also relativ prim zu $2$, und daher ist $\mathcal{E}(1)$ und folglich auch $\mathcal{E}(r)$ durch $(1-r)$ theilbar. Dies aber widerspricht der Voraussetzung, dass $\mathcal{E}(r)$ eine Einheit sein soll.

Demnach ergiebt sich aus (2):
\[
r^{-\frac{1}{2}h}\,\mathcal{E}(r) = r^{\frac{1}{2}h}\,\mathcal{E}(r^{-1}),
\]
woraus folgt, dass $r^{-\frac{1}{2}h}\,\mathcal{E}(r)$ eine reelle Einheit ist. Bezeichnen wir sie mit $e(r)$, so ist also
\begin{equation}\label{eq:178-3}
\mathcal{E}(r) = r^{\frac{1}{2}h}\, e(r),
\end{equation}
woraus für diesen Fall das Theorem 1. bewiesen ist.

Wenn aber zweitens $m$ eine Potenz von $2$ ist, so ist $\mu = \frac{1}{2}m$ und $r^\mu = -1$, und zugleich mit $r$ ist $-r$ eine $m^{\text{te}}$ Einheitswurzel. Wenn wir nöthigenfalls $h$ durch $h + \frac{1}{2}m$ ersetzen, so können wir in (2) das obere Zeichen annehmen und setzen:
\begin{equation}\label{eq:178-4}
\mathcal{E}(r) = r^h\, \mathcal{E}(r^{-1}).
\end{equation}

Es kommt jetzt darauf an, nachzuweisen, dass $h$ gerade sein muss.

Wenn wir $\mathcal{E}(r)$ durch die Basis $1, r, r^2, \ldots, r^{\mu-1}$ darstellen, so ergiebt sich
\begin{equation}\label{eq:178-5}
\mathcal{E}(r) = \sum_{0,\mu-1} x_s\, r^s,
\end{equation}
worin die $x_s$ ganze rationale Zahlen sind. Nach (4) ist aber dann
\[
\sum_{0,\mu-1} x_s\, r^s = \sum_{1,\mu-1} x_s\, r^{h-s},
\]
und daraus folgt, dass $x_s = \pm x_{s'}$ ist, wenn $s + s' \equiv h \pmod{\mu}$. Ist nun $h$ ungerade, so ist aus zwei so verbundenen Zahlen die eine gerade, die andere ungerade (da $\mu$ eine Potenz von $2$ ist). Der Ausdruck (5) ergiebt also für $\mathcal{E}(r)$ ein Aggregat von Gliedern der Form
\[
x_s(r^s \pm r^{s'}).
\]

Es ist aber $r^s \pm r^{s'} = r^s(1 \pm r^{h-2s}) = r^s(1 \mp r^{\mu+h-2s})$ immer durch $1 - r$ theilbar, und daraus würde folgen, dass $\mathcal{E}(r)$ durch $1 - r$, was keine Einheit ist, theilbar wäre, was dem Begriffe der Einheit widerspricht. Also ist die Annahme unzulässig, dass in der Formel (4) der Exponent $h$ ungerade sei, und wir können wie oben
\[
\mathcal{E}(r) = r^{\frac{1}{2}h}\, e(r)
\]
setzen, worin $e(r)$ eine reelle Einheit ist. Damit ist der Satz 1. bewiesen.

Wir wollen ein System von reellen Einheiten hervorheben: Nach \S. 169 sind die $\mu$ Grössen $1 - r^n$ alle unter einander associirt. Demnach ist der Quotient
\[
\frac{1 - r^n}{1 - r^{n'}},
\]
worin $n, n'$ zwei relative Primzahlen zu $m$ bedeuten, eine Einheit. Daraus leitet man aber, wenn
\[
r = e^{\frac{2\pi i}{m}}
\]
gesetzt wird, die reelle Einheit
\begin{equation}\label{eq:178-6}
\frac{\sin\frac{n\pi}{m}}{\sin\frac{n'\pi}{m}} = \frac{r^{\frac{n}{2}} - r^{-\frac{n}{2}}}{r^{\frac{n'}{2}} - r^{-\frac{n'}{2}}} = r^{\frac{n'-n}{2}}\,\frac{1 - r^n}{1 - r^{n'}}
\end{equation}
her. Ist $m$ gerade, so kann man $n' = n + \frac{m}{2}$ setzen, und erhält für diesen Fall die reellen Einheiten
\begin{equation}\label{eq:178-7}
\tau_n = \tang\frac{n\pi}{m},
\end{equation}
worin $n$ jede ungerade Zahl bedeuten kann.

Ersetzt man $n$ durch $\frac{1}{2}m - n$, so geht $\tau_n$ in den reciproken Werth über. Lässt man $n$ die Reihe der Zahlen
\[
1,\; 3,\; 5,\; \ldots,\; \frac{m}{4} - 1
\]
durchlaufen, so erhält $\tau_n$ lauter positive echt gebrochene Werthe.

\bigskip
\begin{center}
\textbf{Zweiundzwanzigster Abschnitt.}\\[4pt]
\textbf{Abel'sche Körper und Kreistheilungskörper.}
\end{center}
\bigskip

\noindent\textbf{\S. 179. Einfache Abel'sche Körper}\\[4pt]

Wir wenden uns nun zu einer der interessantesten Anwendungen der Theorie der algebraischen Zahlen.

Es handelt sich um den Beweis des allgemeinen Satzes\footnote{Kronecker hat diesen Satz zuerst ausgesprochen, aber keinen vollständigen Beweis dafür veröffentlicht. Den ersten Beweis hat der Verfasser des vorliegenden Werkes in Bd. 8 der Acta Mathematica bekannt gemacht (1886). In neuester Zeit ist ein zweiter Beweis von Hilbert veröffentlicht, der sich auf die im neunzehnten Abschnitte entwickelten Sätze stützt (Nachrichten d.\ Gesellschaft d.\ Wissenschaften zu Göttingen, 1896, Heft 1).}:

\begin{satz}
I. Alle im absoluten Rationalitätsbereich Abel'schen Zahlkörper sind Kreistheilungskörper.
\end{satz}

Haben wir diesen Satz bewiesen, so gewinnen die Untersuchungen des dritten Abschnittes über Kreistheilungskörper ein erhöhtes Interesse, weil damit gezeigt ist, dass auf dem dort angegebenen Wege nicht nur alle Kreistheilungskörper, sondern alle Abel'schen Körper überhaupt gefunden werden.

Es sei $\Omega = R(x)$ ein Abel'scher Körper $m^{\text{ten}}$ Grades, d.\,h.\ ein Körper, der aus allen rationalen Functionen einer Wurzel $x$ einer irreduciblen Abel'schen Gleichung $m^{\text{ten}}$ Grades besteht, wenn als Rationalitätsbereich der Körper $R$ der rationalen Zahlen betrachtet wird. Die Galois'sche Gruppe $\mathfrak{G}$ dieser Gleichung, die also eine Abel'sche Gruppe ist und denselben Grad $m$ hat, wie der Körper $\Omega$, ist auch die Gruppe des Körpers $\Omega$.

Ist $x_1$ eine nicht primitive Zahl aus $\Omega$, so ist $\Omega_1 = R(x_1)$ gleichfalls ein Abel'scher Körper, den wir einen echten Theiler von $\Omega$ nennen und dessen Grad ein echter Theiler des Grades von $\Omega$ ist. Denn ist $\mathfrak{A}$ die Gruppe, zu der die Zahl $x_1$ gehört, so ist $\mathfrak{G}/\mathfrak{A}$ (\S. 4 dieses Bandes) die Gruppe des Körpers $\Omega_1$, und dies ist zugleich mit $\mathfrak{G}$ eine Abel'sche Gruppe. Ein Theiler von $\Omega$, der mit $\Omega$ von gleichem Grade ist, ist mit $\Omega$ identisch (vgl. Bd. I, \S. 144, 149, 156).

Giebt es nun zwei nicht primitive Zahlen $x_1, x_2$ in $\Omega$ von der Art, dass
\[
\Omega = R(x_1, x_2)
\]
gesetzt werden kann, so sind die Körper $\Omega_1 = R(x_1)$, $\Omega_2 = R(x_2)$ echte Theiler von $\Omega$, und $\Omega$ heisst aus den beiden Körpern $\Omega_1, \Omega_2$ \emph{zusammengesetzt}.

Gestattet der Körper $\Omega$ keine solche Darstellung, so heisst er \emph{einfach}.

Wenn also als bewiesen vorausgesetzt wird, dass $\Omega_1, \Omega_2$ Kreistheilungskörper sind, d.\,h.\ dass alle ihre Zahlen rational durch Einheitswurzeln darstellbar sind, so folgt dasselbe für $\Omega$.

Wenn einer der Körper $\Omega_1, \Omega_2$ zerlegbar ist, so können wir die Zerlegung wiederholen, und müssen, da die Grade abnehmende ganze positive Zahlen sind, endlich zu einfachen Körpern gelangen.

Das Theorem I. braucht also nur für einfache Abel'sche Körper bewiesen zu werden.

Es ist aber daran zu erinnern, dass die Zerlegbarkeit eines Körpers $\Omega$ keineswegs schon aus der Existenz eines Theilers folgt, und dass bei einem zerlegbaren Körper die Zerlegung in einfache Körper auf ganz verschiedene Arten geschehen kann, so dass bei den Körpern nicht die Gesetze der Theilbarkeit gelten, wie im Gebiete der Zahlen.

Ein Kennzeichen für die Einfachheit eines Körpers können wir aus seiner Gruppe ableiten.

Wenn die Gruppe $\mathfrak{G}$ zwei echte Theiler $\mathfrak{A}$ und $\mathfrak{B}$ von der Art hat, dass jedes Element ein- und nur einmal aus einem Element von $\mathfrak{A}$ und einem Element von $\mathfrak{B}$ zusammengesetzt werden kann, so nennen wir die Gruppe $\mathfrak{G}$ zerlegbar in die beiden Componenten $\mathfrak{A}, \mathfrak{B}$; $\mathfrak{G}$ ist also zerlegbar, wenn in dem nach der Composition der Theile (\S. 4) gebildeten Producte
\begin{equation}\label{eq:179-1}
\mathfrak{G} = \mathfrak{A}\,\mathfrak{B}
\end{equation}
jedes Element von $\mathfrak{G}$ ein- und nur einmal erscheint. Der Grad von $\mathfrak{G}$ ist dann gleich dem Producte der Grade von $\mathfrak{A}$ und $\mathfrak{B}$.

Giebt es solche Theiler nicht, so heisst die Gruppe $\mathfrak{G}$ unzerlegbar.

Dann können wir den Satz beweisen:

\begin{satz}
1. Ist die Gruppe eines Abel'schen Körpers zerlegbar, so ist auch der Körper zusammengesetzt.
\end{satz}

Nehmen wir, um dies zu beweisen, an, die Gruppe $\mathfrak{G}$ des Körpers $\Omega$ sei nach der Formel (1) zerlegbar und $m, a, b$ seien die Grade von $\mathfrak{G}, \mathfrak{A}, \mathfrak{B}$. Wir nehmen eine zur Gruppe $\mathfrak{B}$ gehörige Zahl $\xi$ des Körpers $\Omega$, die also durch die Substitutionen von $\mathfrak{G}$ in $a$ verschiedene conjugirte Werthe übergeht und ebenso eine zu $\mathfrak{A}$ gehörige $b$-wertige Zahl $\eta$. Wir können dann eine Zahl
\[
x = \alpha\,\xi + \beta\,\eta
\]
mit rationalen Zahlencoëfficienten $\alpha, \beta$ bilden, die $m$ verschiedene conjugirte Werthe hat, und dann ist $\Omega = R(x)$. Also ist $\Omega$ aus $R(\xi)$ und $R(\eta)$ zusammengesetzt (Bd. I, \S. 143).

Erinnern wir uns nun an die in \S. 9, 10 dieses Bandes abgeleitete Darstellung einer Abel'schen Gruppe durch eine Basis, so ergiebt sich durch wiederholte Anwendung des eben Bewiesenen:

\begin{satz}
2. Jeder Abel'sche Körper $\Omega$ lässt sich aus solchen zusammensetzen, deren Grad eine Primzahlpotenz und deren Gruppe cyklisch ist. Die Grade dieser zusammensetzenden Körper sind die Invarianten der Gruppe von $\Omega$.
\end{satz}

Diese Sätze werden durch den folgenden ergänzt:

\begin{satz}
3. Ein Abel'scher Körper von Primzahlpotenzgrad mit cyklischer Gruppe ist einfach.
\end{satz}

Wenn die Gruppe $\mathfrak{G}$ cyklisch ist, so lassen sich ihre Elemente durch die Potenzen einer Basis, $G^\gamma$, darstellen, wenn $\gamma$ ein volles Restsystem nach dem Modul $m$ durchläuft. Wir erhalten alle Theiler $\mathfrak{A}$ von $\mathfrak{G}$ dargestellt durch
\[
\mathfrak{A} = G^{b\gamma},
\]
wenn $b$ ein Theiler von $m$ ist und $\gamma$ ein volles Restsystem nach dem Modul $a = m : b$ durchläuft.

Ist
\[
\mathfrak{A}' = G^{b'\gamma}
\]
ein zweiter Theiler von $\mathfrak{G}$ vom Grade $a'$ und ist
\[
b' \gtrless b,\quad a' \lessgtr a,
\]
so ist, wenn wir jetzt annehmen, dass $m$ und mithin auch $a, b, a', b'$ Potenzen einer Primzahl $q$ sind, $b'$ ein Theiler von $b$ und folglich $\mathfrak{A}$ ein Theiler von $\mathfrak{A}'$.

Wenn also $\xi$ eine Zahl des Körpers $\Omega$ ist, die zu der Gruppe $\mathfrak{A}$ gehört, so ist $\xi$ eine primitive Zahl des Körpers $b^{\text{ten}}$ Grades $R(\xi)$, und in diesem Körper ist auch jede Zahl $\xi'$ enthalten, die die Substitutionen der Gruppe $\mathfrak{A}'$ gestattet. Sind also $\mathfrak{A}, \mathfrak{B}$ und $\mathfrak{A}', \mathfrak{B}'$ echte Theiler von $\mathfrak{G}$, so ist $R(\xi)$ von $\Omega$ verschieden, und $R(\xi)$ wird durch Adjunction von $\xi'$ nicht erweitert. Es kann also auch nicht $\Omega$ aus $R(\xi)$ und $R(\xi')$ zusammengesetzt werden, wodurch 3. bewiesen ist.

\bigskip\noindent\textbf{\S. 180. Die Resolventen}\\[4pt]

Nach dem, was jetzt bewiesen ist, können wir uns in den weiteren Betrachtungen, die den Beweis des Satzes I. zum Ziele haben, auf solche Abel'sche Körper beschränken, deren Grad eine Primzahlpotenz und deren Gruppe cyklisch ist. Aus diesem Grunde genügte es auch für die beabsichtigte Anwendung, dass wir im vorigen Abschnitte uns auf die Betrachtung der Kreistheilungskörper $\Omega_m$ beschränkt haben, in denen $m$ eine Primzahlpotenz war. Wir behalten so viel als möglich die dort gebrauchte Bezeichnung bei und setzen
\begin{equation}\label{eq:180-1}
m = q^\varkappa,
\end{equation}
worin $q$ eine Primzahl, $\varkappa$ ein positiver Exponent ist, und in dem Falle, dass $q = 2$ ist, nehmen wir $\varkappa > 1$ an. Mit $n$ bezeichnen wir jede durch $q$ nicht theilbare Zahl. Es ist $r$ eine primitive $m^{\text{te}}$ Einheitswurzel und $\Omega_m$ der Körper der rationalen Functionen von $r$.

Es sei nun $\Omega = R(x)$ ein einfacher Abel'scher Körper $m^{\text{ten}}$ Grades, so dass $x$ Wurzel einer rationalen irreduciblen cyklischen Gleichung $m^{\text{ten}}$ Grades ist. Bezeichnen wir die Wurzeln dieser Gleichung in geeigneter Reihenfolge mit
\begin{equation}\label{eq:180-2}
x_0,\; x_1,\; x_2,\; \ldots,\; x_{m-1},
\end{equation}
und nehmen $x_m = x_0$, $x_{m+1} = x_1$, \ldots\ an, so ist
\[
x_1 = \Phi(x_0),\; x_2 = \Phi(x_1),\; \ldots,\; x_0 = \Phi(x_{m-1}),
\]
worin $\Phi$ eine ganze Function mit rationalen Coëfficienten bedeutet, und die Gruppe des Körpers $\Omega$ besteht aus den Potenzen der Substitution $(x_0, x_1)$ oder aus den cyklischen Permutationen der Wurzeln (2). Jede cyklische Function dieser Grössen ist eine rationale Zahl. Was wir zu beweisen haben, ist, dass sich die $x$ rational durch Einheitswurzeln ausdrücken lassen. Adjungiren wir dem Körper $\Omega$ noch die $m^{\text{te}}$ Einheitswurzel $r$, so entsteht ein Körper $R(x, r)$, der die beiden Körper $\Omega = R(x)$ und $R(r) = \Omega_m$ als Theiler enthält.

Wenn eine Zahl $\omega = \Phi(x, r)$ des Körpers $R(x, r)$ die sämmtlichen Substitutionen $(x_0, x_1)$, $(x_0, x_2)$, \ldots, $(x_0, x_{m-1})$ gestattet, so ist sie im Körper $\Omega_m$ enthalten; denn dann ist
\[
m\omega = \Phi(x_0, r) + \Phi(x_1, r) + \cdots + \Phi(x_{m-1}, r),
\]
und darin sind die Coëfficienten der Potenzen von $r$ nicht nur cyklische, sondern sogar symmetrische Functionen der Wurzeln (2), also rationale Zahlen. Dies gilt selbst noch in dem Falle, dass die Gleichung für $x$ durch Adjunction von $r$ reducirt wird.

Dies wenden wir an auf die Resolventen
\begin{equation}\label{eq:180-3}
(r, x_0) = x_0 + r x_1 + r^2 x_2 + \cdots + r^{m-1} x_{m-1},
\end{equation}
und allgemeiner, wenn $s$ einen beliebigen ganzzahligen Exponenten bedeutet:
\begin{equation}\label{eq:180-4}
(r^s, x_0) = x_0 + r^s x_1 + r^{2s} x_2 + \cdots + r^{(m-1)s} x_{m-1}.
\end{equation}

Durch diese Resolventen lässt sich $x_0$ selbst wieder rational ausdrücken, nämlich:
\begin{equation}\label{eq:180-5}
m x_0 = \sum_{0,m-1} (r^s, x_0),
\end{equation}
und wenn wir also nachweisen können, dass alle diese Resolventen $(r^s, x_0)$ Kreistheilungszahlen sind, so haben wir unser Ziel erreicht. Wenn wir
\begin{equation}\label{eq:180-6}
(r^s, x_\varkappa) = x_\varkappa + r^s x_{\varkappa+1} + r^{2s} x_{\varkappa+2} + \cdots + r^{(m-1)s} x_{\varkappa+m-1}
\end{equation}
setzen, so geht $(r^s, x_\varkappa)$ aus $(r^s, x_0)$ durch die Substitution $(x_0, x_\varkappa)$ hervor, und es besteht die fundamentale Relation
\begin{equation}\label{eq:180-7}
r^{\varkappa s}(r^s, x_\varkappa) = (r^s, x_0).
\end{equation}

Alle diese Resolventen sind Zahlen des Körpers $R(x, r)$; wir betrachten vorzugsweise die beiden folgenden Verbindungen:
\begin{align}
(r^s, x_0)^m &= \omega_s,\label{eq:180-8}\\
(r^s, x_0)^{-1}(r, x_0)^s &= \alpha.\label{eq:180-9}
\end{align}

\bigskip\noindent\textbf{\S. 181. Vorbereitung zum Beweis}\\[4pt]

Nach den Ausführungen des vorigen Paragraphen ist zum Beweis des grossen Theorems I. nur noch der Nachweis erforderlich, dass die in $\Omega_m$ enthaltene Zahl
\[
\omega = (r, x_0)^m
\]
die $m^{\text{te}}$ Potenz einer Kreistheilungszahl ist.

\begin{satz}
1. Die wesentliche Eigenschaft der zu $\omega$ conjugirten Zahlen $\omega_n$, auf die sich der Beweis stützt, ist die, dass
\begin{equation}\label{eq:181-1}
\omega_n^{-1}\,\omega_1^n = \alpha^m
\end{equation}
die $m^{\text{te}}$ Potenz einer Zahl in $\Omega_m$ ist.
\end{satz}

Hierin kann $n$ jede durch $q$ nicht theilbare Zahl sein. Bilden wir von der rechten und linken Seite von (1) die Normen im Körper $\Omega_m$, so ergiebt sich, wenn wir die Norm der Zahl $\alpha$, die eine rationale Zahl ist, mit $a$ bezeichnen, da jedes $\omega_n$ dieselbe Norm wie $\omega$ hat,
\[
N(\omega)^{n-1} = a^m.
\]

Wenn nun $m$ ungerade ist, so können wir hierin $n = 2$ annehmen, und erhalten als Folgerung aus 1.:

\begin{satz}
2. Die Norm von $\omega$ ist die $m^{\text{te}}$ Potenz einer rationalen Zahl
\begin{equation}\label{eq:181-2}
N(\omega) = a^m.
\end{equation}
\end{satz}

Ist $m$ eine Potenz von 2, so ist diese Folgerung nicht mehr zulässig. Es ergiebt sich dann aus 1. nur, dass $N(\omega)^2$ die $m^{\text{te}}$ Potenz einer rationalen Zahl ist. Gleichwohl müssen die Zahlen $\omega$, wie aus (10) des vorigen Paragraphen zu sehen ist, auch in diesem Falle noch der Bedingung 2. genügen, die dann aber nicht mehr von selbst schon in 1. enthalten ist, sondern als neue Eigenschaft hinzukommt.

Was wir zu beweisen haben, ist, dass aus 1. und 2. folgt, dass auch $\omega$ selbst die $m^{\text{te}}$ Potenz einer Kreistheilungszahl, wenn auch in einem höheren Körper, ist.

Dazu ist aber erforderlich, die Zerlegung der Zahlen $\omega$ in ihre Primfactoren im Körper $\Omega_m$ zu untersuchen.

Die Primzahl $q$ ist, wie im \S. 169 nachgewiesen, mit der $\varphi(m)^{\text{ten}}$ Potenz einer in $\Omega_m$ existirenden Primzahl $\sigma = 1 - r$ associirt. Ist also $\sigma^k$ die in $\omega$ enthaltene Potenz von $\sigma$, so setzen wir
\[
\omega = \sigma^k \omega',
\]
worin $\omega'$ mit der Primzahl $q$ theilerfremd ist. Die Bildung der Norm ergiebt nach 2.
\[
a^m = q^k b,
\]
worin $b$ eine rationale Zahl ist, die in einfachster Gestalt die Primzahl $q$ weder im Zähler noch im Nenner enthält. Daraus ergiebt sich aber, dass $k$ durch $m$ theilbar sein muss, und daraus das erste Resultat:

\begin{satz}
3. Der Exponent der in $\omega$ enthaltenen Primzahl $\sigma$ ist immer durch $m$ theilbar.
\end{satz}

Es sei nun ferner $p$ eine von $q$ verschiedene Primzahl und $p_\xi$ ein Primfactor von $p$ im Körper $\Omega_m$. Wir lassen $\xi$ die im \S. 172 in den verschiedenen Fällen näher bestimmte Zahlenreihe $\xi_1, \xi_2, \ldots, \xi_e$ durchlaufen, so dass
\begin{equation}\label{eq:181-3}
p = p_{\xi_1} p_{\xi_2} \ldots p_{\xi_e}
\end{equation}
wird, und nehmen an, es sei $p_\xi^{k_\xi}$ die in $\omega$ enthaltene Potenz von $p_\xi$. Die in $\omega_n$ enthaltene Potenz von $p_{n\xi}$ ist dann $p_{n\xi}^{k_\xi}$, und wenn wir $n'$ aus der Congruenz
\begin{equation}\label{eq:181-4}
nn' \equiv 1 \pmod{m}
\end{equation}
bestimmen, so ist $p_\xi^{k_{n'\xi}}$ die in $\omega_n$ enthaltene Potenz von $p_\xi$. Demnach ist in
\begin{equation}\label{eq:181-5}
\omega_n^{-1}\,\omega_1^n \text{ die Potenz } p_\xi^{nk_\xi - k_{n'\xi}}
\end{equation}
enthalten, und wegen 1. muss der Exponent dieser Potenz durch $m$ theilbar sein.

Wir erhalten also für jedes durch $q$ nicht theilbare $n$ die Congruenz
\begin{equation}\label{eq:181-6}
nk_\xi \equiv k_{n'\xi} \pmod{m},
\end{equation}
in der wir nun auch $\xi$ durch $n\xi$ ersetzen können, wodurch sich
\begin{equation}\label{eq:181-7}
nk_{n\xi} \equiv k_\xi \pmod{m}
\end{equation}
ergiebt. Nehmen wir hierin $\xi = 1$ und setzen dann $\xi, \xi'$ für $n, n'$ und $k$ für $k_1$, so folgt aus (7)
\begin{equation}\label{eq:181-8}
k_\xi \equiv \xi' k \pmod{m},
\end{equation}
worin nun $k$ für alle $\xi$ denselben Werth hat und $\xi, \xi'$ durch die Congruenz
\[
\xi\xi' \equiv 1 \pmod{m}
\]
zusammenhängen.

Nehmen wir $n = p$ an, so wird $p_{p\xi} = p_\xi$ (\S. 172), also ist auch $k_{p\xi} = k_\xi$, und aus (7) und (8) folgt:
\begin{equation}\label{eq:181-9}
k(p-1) \equiv 0 \pmod{m}.
\end{equation}

Hieraus ergeben sich nun wichtige Folgerungen: Wenn zunächst $p-1$ nicht durch $q$ theilbar ist, so folgt, dass $k$ durch $m$ theilbar sein muss, und wir haben also den Satz:

\begin{satz}
4. Die Primfactoren einer Primzahl $p$, die nach dem Modul $q$ nicht mit $1$ congruent ist, sind in $\omega$ nur in solchen Potenzen enthalten, deren Exponenten durch $m$ theilbar sind.
\end{satz}

Wenn in dem Falle, wo $m$ eine Potenz von $2$ ist, $p-1$ durch $2$, aber nicht durch $4$ theilbar ist, wenn also $p \equiv -1 \pmod{4}$ ist, so folgt aus (9) nur, dass $k$ durch $\frac{1}{2}m$, nicht aber, dass $k$ durch $m$ theilbar ist. Der Exponent der in $\omega$ enthaltenen Potenz von $p_\xi$ ist, von Vielfachen von $m$ abgesehen, gleich $k$, und wenn $k$ durch $m$ theilbar ist, so ist alles wie im Falle 4. Ist aber $k$ nur durch $\frac{1}{2}m$, nicht durch $m$ theilbar, also
\[
k \equiv \tfrac{1}{2}m \pmod{m},
\]
so ist $k - \frac{1}{2}m$ durch $m$ theilbar, und wir können mit Rücksicht auf die Zerlegung (3) den Satz so formuliren:

\begin{satz}
5. Ist $m$ eine Potenz von $2$ und $p$ eine Primzahl congruent mit $-1$ nach dem Modul $4$, so lässt sich der Exponent $h$ so bestimmen, dass alle Primfactoren von $p$ in $\omega\sqrt[m]{p}^{-hm}$ nur zu solchen Potenzen enthalten sind, deren Exponent durch $m$ theilbar ist.
\end{satz}

Es sei endlich $p-1$ durch $q$, oder, falls $m$ gerade ist, durch $4$ theilbar. Ist $m_1$ der grösste gemeinschaftliche Theiler von $m$ und $p-1$, und ist
\[
m = m_1 m_2,
\]
so muss $k$ nach (9) durch $m_2$ theilbar sein, und wenn wir
\[
k = h m_2
\]
setzen, so ist in $\omega$ nach (8) das Product enthalten:
\begin{equation}\label{eq:181-10}
\bigl(\prod^\xi p_\xi^{\xi'}\bigr)^{hm_2}.
\end{equation}
Ausserdem kommen in $\omega$ die Primfactoren $p_\xi$ nur noch in solchen Potenzen vor, deren Exponenten durch $m$ theilbar sind.

In diesem Falle durchläuft nun $\xi$ nach \S. 172, III ein volles System durch $q$ nicht theilbarer Reste nach dem Modul $m_1$, und wir können in (10) unter $\xi'$ die kleinste positive Lösung der Congruenz
\[
\xi\xi' \equiv 1 \pmod{m_1}
\]
verstehen, weil, wenn wir die Exponenten $\xi'$ in (10) nach dem Modul $m_1$ verändern, nur $m^{\text{te}}$ Potenzen der Primfactoren $p$ hinzutreten.

Dann können wir auf das Product (10) das Kummer'sche Theorem [\S. 174, (16)] anwenden, indem wir dort $m$ durch $m_1$ ersetzen, und erhalten, wenn wir unter $r_1 = r^{m_2}$ eine $m_1^{\text{te}}$ Einheitswurzel, unter den $\eta$ gewisse Perioden der $p^{\text{ten}}$ Einheitswurzeln verstehen:
\[
\prod^\xi p_\xi^{\xi'} = (r_1, \eta)^{m_1},
\]
und folglich
\begin{equation}\label{eq:181-11}
\bigl(\prod^\xi p_\xi^{\xi'}\bigr)^{hm_2} = (r_1, \eta)^{hm}.
\end{equation}

Demnach haben wir das Theorem:

\begin{satz}
6. Ist $p$ eine Primzahl und $p-1$ durch $q$, oder, wenn $q = 2$ ist, durch $4$ theilbar, so lässt sich der positive Exponent $h$ so bestimmen, dass die in $\Omega_m$ enthaltene Zahl $\omega\,(r_1, \eta)^{-hm}$ die Primfactoren von $p$ nur zu solchen Potenzen enthält, deren Exponent durch $m$ theilbar ist.
\end{satz}

Es ist jetzt auch nicht mehr nöthig, den Satz 5. von dem Satze 6. zu unterscheiden; denn für $m_1 = 2$ geht 6. in 5. über, weil dann $r_1 = -1$ wird, und nach Bd. I, \S. 171
\[
(r_1, \eta)^m = (-1, \eta)^m = \sqrt[m]{p}^m
\]
ist\footnote{In des Verfassers Abhandlung in Bd. 8 der Acta mathematica ist es versäumt, den Fall des Satzes 5. besonders hervorzuheben.}.

Fassen wir das Ergebniss der Sätze 3. bis 6. in eine Formel zusammen, so ergiebt sich, wenn $\varepsilon$ ein Einheitsfunctional, $\varphi$ irgend ein Functional des Körpers $\Omega_m$ bedeutet, und $p$ eine Reihe von Primzahlen durchläuft, die congruent mit $1$ nach dem Modul $q$ sind, worunter dieselbe Primzahl $p$ auch mehrmals vorkommen kann,
\begin{equation}\label{eq:181-12}
\omega = (r, x_0)^m = \varepsilon\,\varphi^m \prod^p (r^{m_2}, \eta)^m.
\end{equation}

Die in diesen Formeln vorkommenden Resolventen der Kreistheilung $(r^{m_2}, \eta)$ sind Kreistheilungszahlen in einem höheren Körper. Ihre $m^{\text{ten}}$ Potenzen sind aber im Körper $\Omega_m$ enthaltene Zahlen, die durch die Substitution $(r, r^m)$ in $(r^{mm_2}, \eta)^m$ übergehen. Ebenso sind die Verbindungen
\begin{equation}\label{eq:181-13}
(r^{mm_2}, \eta)^{-1}(r^{m_2}, \eta)^n
\end{equation}
im Körper $\Omega_m$ enthalten [\S. 174, (9), (11)]. Diese Resolventen sind specielle Fälle der allgemeinen Resolvente $(r, x_0)$. Aus (12) erhalten wir, wenn wir mit $\varphi_n$, $\varepsilon_n$ die conjugirten Functionale zu $\varphi$ und $\varepsilon$ bezeichnen
\begin{equation}\label{eq:181-14}
\omega_n = \varepsilon_n\,\varphi_n^m \prod^p (r^{m_2 n}, \eta)^m.
\end{equation}

Diese Formel lehrt uns die wichtigsten Eigenschaften der Functionale $\varphi_n$ und $\varepsilon_n$ kennen, zunächst:

\begin{satz}
7. Die Functionale $\varphi_n^m$ sind mit Zahlen des Körpers $\Omega_m$ associirt, gehören also der Hauptclasse an (\S. 152).
\end{satz}

Nach 1. ist $\omega_n^{-1}\,\omega_1^n = \alpha^m$ die $m^{\text{te}}$ Potenz einer Zahl $\Omega_m$. Da auch die Verbindungen (13) Zahlen in $\Omega_m$ sind, so können wir
\[
\prod^p (r^{m_2 n}, \eta)^{-1}(r^{m_2}, \eta)^n = \frac{\alpha}{\beta}
\]
setzen, und erhalten aus (14)
\[
\varepsilon_n^{-1}\,\varepsilon_1^n\,(\varphi_n^{-1}\,\varphi_1^n)^m = \beta^m,
\]
worin $\beta$ eine Zahl in $\Omega_m$ ist.

Daraus geht hervor, dass das Product
\[
\beta\,\varphi_n\,\varphi_1^{-n} = \varepsilon
\]
ein Einheitsfunctional des Körpers $\Omega_m$ ist, und daraus ergeben sich für die Functionale $\varphi_n$ und die Einheitsfunctionale $\varepsilon_n$ die folgenden beiden Sätze:

\begin{satz}
8. Die conjugirten Functionale $\varphi_n$ haben die Eigenschaft, dass alle Producte $\varphi_n\,\varphi_1^{-n}$ der Hauptclasse angehören.
\end{satz}

\begin{satz}
9. Die Einheitsfunctionale $\varepsilon_n$ haben die Eigenschaft, dass die Producte $\varepsilon_n\,\varepsilon_1^{-n} = \varepsilon^m$ $m^{\text{te}}$ Potenzen von Einheitsfunctionalen sind.
\end{satz}

\bigskip\noindent\textbf{\S. 182. Beweis des ersten Hülfsatzes für ein ungerades $m$}\\[4pt]

Es kommt nun für den Beweis des Haupttheorems I. alles auf den Beweis der beiden Sätze an, die aus den Eigenschaften der Functionale $\varphi$ und $\varepsilon$ zu folgern sind:

\textbf{A.} Die Functionale $\varphi_n$ gehören der Hauptclasse an.

Können wir dieses Lemma voraussetzen, so ist $\varphi_n$ mit einer Zahl $\alpha_n$ associirt, und wir können die Formel (14) mit etwas veränderter Bedeutung von $\varepsilon_n$ so darstellen:
\begin{equation}\label{eq:182-15}
\omega_n = \varepsilon_n\,\alpha_n^m \prod^p (r^{m_2}, \eta)^m.
\end{equation}

In dieser Formel ist aber $\varepsilon_n$ eine numerische Einheit des Körpers $\Omega_m$, die dem Satze 9. genügt.

Können wir also aus dem Satze 9. noch das zweite Lemma beweisen:

\textbf{B.} Eine numerische Einheit $\varepsilon_n$ des Körpers $\Omega_m$, die dem Satze 9. genügt, ist die $m^{\text{te}}$ Potenz einer Einheit in $\Omega_m$, multiplicirt mit einer Potenz von $r$;

so können wir in (15) die $m^{\text{te}}$ Wurzel ziehen, und erhalten, wenn mit $\varrho$ eine Einheitswurzel von möglicherweise höherem Grade als $m$ und mit $\alpha$ eine Zahl des Körpers $\Omega_m$ bezeichnet wird:
\begin{equation}\label{eq:182-16}
(r, x_0) = \varrho\,\alpha \prod^p (r^{m_2}, \eta),
\end{equation}
wo nun rechts lauter Kreistheilungszahlen stehen und wodurch daher das Theorem I. erwiesen ist.

Wir haben im vorigen Paragraphen den Beweis des Theorems I. von zwei Hülfssätzen abhängig gemacht, deren Beweis nun ferner zu suchen ist.

Wir formuliren den ersten dieser Hülfssätze so:

\textbf{a)} Ist $\varphi_n$ ein System von Null verschiedener conjugirter Functionale des Körpers $\Omega_m$, von denen bekannt ist, dass $\varphi_n^m$ und $\varphi_n^{-1}\,\varphi_1^n$ der Hauptclasse angehören, so gehören die Functionale $\varphi$ selbst der Hauptclasse an.

Der Satz wird bewiesen sein, wenn von irgend einer Potenz $\varphi^k$ von $\varphi$, deren Exponent nicht durch $q$ theilbar ist, bewiesen werden kann, dass sie der Hauptclasse angehört. Denn sind $\alpha, \beta$ Zahlen in $\Omega_m$ und ist
\[
\varphi^m = \varepsilon'\alpha,\quad \varphi^k = \varepsilon''\beta,
\]
so können wir die ganzen rationalen Zahlen $x, y$ so bestimmen, dass
\[
mx + ky = 1
\]
wird, und dann wird
\[
\varphi = \varphi^{mx}\,\varphi^{ky} = \varepsilon'''\alpha^x\beta^y,
\]
also ist auch, wenn $\varepsilon', \varepsilon'', \varepsilon'''$ Einheiten sind, $\varphi$ selbst mit einer Zahl associirt.

Nach dem heutigen Stande der Sache ist der Beweis für ein gerades $m$ auf einem ganz anderen Wege zu führen, als für ein ungerades, und wir beschäftigen uns also zunächst mit dem leichteren Falle des ungeraden $m$\footnote{Vgl. übrigens den neuen Beweis von Hilbert, bei dem der Unterschied zwischen geradem und ungeradem $m$ weit mehr zurücktritt (S. 648).}.

Die in a) über $\varphi_n$ gemachten Voraussetzungen können wir, wenn $\alpha, \beta$ wieder irgend welche Zahlen des Körpers $\Omega_m$ und $\varepsilon$ Einheitsfunctionale bedeuten, durch die beiden Formeln ausdrücken:
\begin{align}
\varphi_n^m &= \varepsilon\,\beta,\label{eq:182-1}\\
\varphi_n &= \varepsilon\,\alpha\,\varphi_1^n.\label{eq:182-2}
\end{align}

Wir werden nun die Aufgabe schrittweise lösen.

Ist zunächst, wie früher, $\sigma$ der in $q$ enthaltene Primfactor, der eine in $\Omega_m$ existirende Zahl ist, und $\sigma^k$ die in $\varphi$ enthaltene Potenz von $\sigma$, so setzen wir
\begin{equation}\label{eq:182-3}
\varphi = \sigma^k \psi,
\end{equation}
worin $\psi$ ein zu $q$ theilerfremdes Functional bedeutet, das denselben Bedingungen (1), (2) wie $\varphi$ genügt, und wenn eines von beiden in die Hauptclasse gehört, so gilt das auch vom anderen. Der Satz a) braucht also nur noch unter der Voraussetzung bewiesen zu werden, dass $\varphi$ theilerfremd zu $q$ ist.

Es sei nun $p$ eine von $q$ verschiedene Primzahl, die in Primfactoren zerlegt den Ausdruck hat:
\begin{equation}\label{eq:182-4}
p = \prod^\xi p_\xi,
\end{equation}
worin $\xi$ die in \S. 172 näher bestimmte Zahlenreihe durchläuft. Es möge $p_\xi^{k_\xi}$ die in $\varphi$ enthaltene Potenz von $p_\xi$ sein, so dass, wenn wir
\[
\varphi = \chi \prod^\xi p_\xi^{k_\xi}
\]
setzen, $\chi$ theilerfremd zu $p$ und mit keinen anderen Primzahlen verwandt ist (s.\ den Schluss von \S. 173), als $\varphi$ selbst. Daraus ergiebt sich
\begin{equation}\label{eq:182-5}
\varphi_n = \chi_n \prod^\xi p_{n\xi}^{k_\xi},
\end{equation}
und hierin lassen wir $n$ abermals die Reihe der Zahlen $\xi$ durchlaufen. Setzen wir noch
\[
\psi = \prod^\xi \chi_\xi,
\]
so ist auch $\psi$ theilerfremd zu $p$ und mit keinen anderen Primzahlen verwandt als $\varphi$, und wir erhalten aus (5) mit Benutzung von (4):
\begin{equation}\label{eq:182-6}
\prod^\xi \varphi_\xi = \psi\,p^{\Sigma k_\xi},
\end{equation}
und daraus ergiebt sich, dass das Functional $\psi$ denselben beiden Bedingungen (1), (2) genügt, wie $\varphi$.

Es zeigt sich ferner, dass, wenn die $\varphi_\xi$ in der Hauptclasse liegen, auch $\psi$ in der Hauptclasse enthalten ist.

Wenn aber $p-1$ nicht durch $q$ theilbar ist, so können wir auch das Umgekehrte schliessen.

Denn nach (2) ist $\varphi_\xi$ äquivalent mit $\varphi^\xi$ und folglich ist:
\begin{align*}
\prod^\xi \varphi_\xi &\text{ äquivalent } \varphi^{\Sigma k_\xi},\\
&\text{äquivalent } \psi,
\end{align*}
und wenn $\psi$ mit einer Zahl associirt ist, so gilt dasselbe von $\varphi^{\Sigma k_\xi}$. Nach dem Satze \S. 172 ist aber, wenn $p-1$ nicht durch $q$ theilbar ist, $\Sigma\,\xi$ auch nicht durch $q$ theilbar, und folglich ist auch $\varphi$ selbst in der Hauptclasse enthalten.

Diese Betrachtung können wir nun, wenn $\psi$ noch mit weiteren Primzahlen, die nach dem Modul $q$ nicht mit $1$ congruent sind, verwandt ist, wiederholen, und kommen so zu dem Satze:

\begin{quote}
Der Satz a) ist allgemein bewiesen, wenn er unter der Voraussetzung bewiesen werden kann, dass $\varphi$ nur mit solchen Primzahlen verwandt ist, die nach dem Modul $q$ mit $1$ congruent sind.
\end{quote}

Machen wir diese Voraussetzung über die Functionale $\varphi_n$, so können wir das Kummer'sche Theorem in der Fassung des \S. 174, 3. anwenden.

Wir bezeichnen in der Folge mit $\varepsilon$ Einheitsfunctionale, auf deren genauere Kenntniss nichts ankommt, und setzen
\begin{equation}\label{eq:182-7}
\Phi_n = \prod^t \varphi_{nt'}^t = \varepsilon\,\vartheta_n^m,
\end{equation}
wenn $t$ die durch $q$ nicht theilbaren Reste von $m$, die kleiner als $m$ sind, durchläuft, und $t'$ durch $tt' \equiv 1 \pmod{m}$ bestimmt ist. Dann ist $\vartheta_n$ eine Kreistheilungszahl in einem höheren Körper, deren $m^{\text{te}}$ Potenz in $\Omega_m$ enthalten ist und die der zweiten Bedingung genügt, dass
\begin{equation}\label{eq:182-8}
\vartheta_n^{-1}\,\vartheta_1^n = \alpha
\end{equation}
eine Zahl des Körpers $\Omega_m$ ist.

Wir führen nun ein vielfach gebrauchtes Zeichen ein, indem wir unter $E(x)$ die grösste ganze Zahl verstehen, die in der reellen Zahl $x$ enthalten ist, und unter $(x)$ den Rest, der stets ein echter Bruch ist, und, wenn $x$ selbst eine ganze Zahl sein sollte, gleich Null zu setzen ist. Dann ist
\begin{equation}\label{eq:182-9}
x = E(x) + (x).
\end{equation}

Ist $x$ eine ganze Zahl, so ist nach dieser Bezeichnungsweise
\[
m\Bigl(\frac{x}{m}\Bigr)
\]
auch eine ganze Zahl, und zwar der kleinste positive Rest der Theilung von $x$ durch $m$.

In der Formel (7) ist der Index der Functionale $\varphi$ nur nach dem Modul $m$ bestimmt. Der Exponent aber ist auf seinen kleinsten Rest nach dem Modul $m$ reducirt. Ersetzen wir also in (7) den Index $t'$ durch $n't'$, so muss $t$ durch den Rest der Division von $nt$ durch $m$, d.\,h.\ [nach (9)] durch
\[
m\Bigl(\frac{nt}{m}\Bigr) = nt - m\,E\Bigl(\frac{nt}{m}\Bigr)
\]
ersetzt werden, und wir können setzen
\[
\Phi_n = \prod^t \varphi_{t'}^{nt - mE(\frac{nt}{m})} = \varepsilon\,\vartheta_n^m,
\]
andererseits ist
\[
\Phi_1 = \prod^t \varphi_{t'}^t = \varepsilon\,\vartheta_1^m,
\]
also nach (8)
\begin{equation}\label{eq:182-10}
\Phi_1^n\,\Phi_n^{-1} = \prod^t \varphi_{t'}^{mE(\frac{nt}{m})} = \varepsilon\,\alpha^m.
\end{equation}

Hiernach ist der Quotient
\[
\Biggl(\frac{\prod^t \varphi_{t'}^{E(\frac{nt}{m})}}{\alpha}\Biggr)^m
\]
ein Einheitsfunctional, und da er zugleich die $m^{\text{te}}$ Potenz eines Functionals in $\Omega_m$ ist, so ist die Basis dieser Potenz auch eine Einheit, d.\,h.\ es ist
\[
\prod^t \varphi_{t'}^{E(\frac{nt}{m})} = \varepsilon\,\alpha,
\]
und es gehört dies Product der Hauptclasse an. Nun ist aber nach unserer Voraussetzung
\[
\varphi_{t'} \text{ äquivalent mit } \varphi^{t'},
\]
und daher ist
\[
\prod^t \varphi_{t'}^{E(\frac{nt}{m})} \text{ äquivalent mit } \varphi^{\sum t'E(\frac{nt}{m})},
\]
und dies gilt für jedes beliebige durch $q$ nicht theilbare $n$.

Können wir also nachweisen, dass sich $n$ so bestimmen lässt, dass auch die Summe
\[
S_n = \sum^t t'\,E\Bigl(\frac{nt}{m}\Bigr)
\]
durch $q$ nicht theilbar ist, so haben wir das Theorem a) bewiesen.

Wir wollen zunächst eine Umformung der Summe $S_n$ vornehmen, durch die die Frage auf den weit einfacheren Fall zurückgeführt wird, in dem $m$ eine Primzahl ist.

Wir setzen, wenn $m$ eine höhere als die erste Potenz von $q$ ist,
\[
m = qm',\quad t = t_1 + qt_2,\qquad \begin{aligned} t_1 &= 1, 2, \ldots, q-1\\[-2pt] t_2 &= 0, 1, \ldots, m'-1.\end{aligned}
\]

Darin ist $m'$ noch eine Potenz von $q$, der Summationsbuchstabe $t_1$ durchläuft die Zahlenreihe $1, 2, \ldots, q-1$, und $t_2$ die Zahlenreihe $0, 1, 2, \ldots, m'-1$. Wir bestimmen $t_1'$ aus der Congruenz
\[
t_1 t_1' \equiv 1 \pmod{q}
\]
und erhalten
\[
t' \equiv t_1' \pmod{q}.
\]

Es ist dann
\begin{equation}\label{eq:182-11}
S_n = \sum^t t'\,E\Bigl(\frac{tn}{m}\Bigr) \equiv \sum^{t_1} t_1' \sum^{t_2} E\Bigl(\frac{nt_2}{m'} + \frac{nt_1}{m}\Bigr) \pmod{q}.
\end{equation}

Wir nehmen jetzt $n < q$ an, also auch $nt_1 < m$, so dass $nt_1 : m$ ein echter Bruch ist. Dann ist
\begin{alignat*}{2}
\text{a)}\qquad &E\Bigl(\frac{nt_2}{m'} + \frac{nt_1}{m}\Bigr) \text{ entweder } &&= E\Bigl(\frac{nt_2}{m'}\Bigr),\\
\text{b)}\qquad & \text{oder } &&= E\Bigl(\frac{nt_2}{m'}\Bigr) + 1,
\end{alignat*}
und zwar tritt der Fall b) nur dann ein, wenn zwischen
\[
\frac{nt_2}{m'} \quad \text{und} \quad \frac{nt_2}{m'} + \frac{nt_1}{m}
\]
eine ganze Zahl liegt. Dies ist aber nur dann der Fall, wenn der Unterschied zwischen $nt_2 : m'$ und der nächst grösseren ganzen Zahl kleiner als $nt_1 : m$ ist, also wenn
\begin{equation}\label{eq:182-12}
E\Bigl(\frac{nt_2}{m'}\Bigr) + 1 - \frac{nt_2}{m'} < \frac{nt_1}{m}.
\end{equation}

Lassen wir in dem Ausdrucke auf der linken Seite von (12)
\begin{equation}\label{eq:182-13}
E\Bigl(\frac{nt_2}{m'}\Bigr) + 1 - \frac{nt_2}{m'}
\end{equation}
$t_2$ alle seine Werthe durchlaufen, so erhalten wir lauter positive, die Einheit nicht übersteigende rationale Brüche mit dem Nenner $m'$, von denen keine zwei einander gleich sind. Denn wären zwei der Werthe, die aus (13) durch Substitution von $t_2', t_2''$ für $t_2$ entstehen, einander gleich, so müsste $\frac{n(t_2' - t_2'')}{m'}$ eine ganze Zahl sein, was, da $n$ relativ prim zu $m'$ und $t_2' - t_2''$ kleiner als $m'$ ist, nicht sein kann. Die Ausdrücke (13) durchlaufen daher in irgend einer Reihenfolge die Zahlenwerthe
\begin{equation}\label{eq:182-14}
\frac{1}{m'},\; \frac{2}{m'},\; \ldots,\; \frac{m'-1}{m'},\; 1,
\end{equation}
und wenn $a$ einen beliebigen positiven Bruch bedeutet, so ist $E(m'a)$ die Anzahl der Zahlen der Reihe (14), die nicht grösser als $a$ sind. Nach (12) ist daher die Anzahl der Werthe von $t_2$, für die der Fall b) eintritt, gleich $E\bigl(\frac{nt_1}{q}\bigr)$, und es ergiebt sich
\[
\sum^{t_2} E\Bigl(\frac{nt_2}{m'} + \frac{nt_1}{m}\Bigr) = \sum^{t_2} E\Bigl(\frac{nt_2}{m'}\Bigr) + E\Bigl(\frac{nt_1}{q}\Bigr).
\]

Um $S_n$ zu bilden, multipliciren wir diese Gleichung mit $t_1'$ und summiren. Dadurch ergiebt sich
\[
S_n \equiv \sum^{t_1} t_1' \sum^{t_2} E\Bigl(\frac{nt_2}{m'}\Bigr) + \sum^{t_1} t_1'\,E\Bigl(\frac{nt_1}{q}\Bigr) \pmod{q}.
\]

Darin ist nun $\sum t_1' = \sum t_1 \equiv 0 \pmod{q}$, weil die $t_1$ paarweise die Summe $q$ ergeben, und es folgt:
\begin{equation}\label{eq:182-15}
S_n \equiv \sum^{t_1} t_1'\,E\Bigl(\frac{nt_1}{q}\Bigr) \pmod{q}.
\end{equation}

Die Summe, die hier auf der rechten Seite steht, ist ebenso gebildet, wie $S_n$ selbst in dem Falle, wo $m$ eine Primzahl ist. Für diesen Fall ergiebt sich aber
\begin{align*}
\frac{nt_1}{q} &= E\Bigl(\frac{nt_1}{q}\Bigr) + \Bigl(\frac{nt_1}{q}\Bigr),\\
\frac{(q-n)t_1}{q} &= E\Bigl(\frac{(q-n)t_1}{q}\Bigr) + \Bigl(\frac{(q-n)t_1}{q}\Bigr),
\end{align*}
und daraus durch Addition
\[
t_1 = E\Bigl(\frac{nt_1}{q}\Bigr) + E\Bigl(\frac{(q-n)t_1}{q}\Bigr) + \Bigl(\frac{nt_1}{q}\Bigr) + \Bigl(\frac{(q-n)t_1}{q}\Bigr),
\]
und dabei kann die Summe der beiden positiven echten Brüche, die doch eine ganze Zahl sein muss, nur den Werth $1$ haben. Es folgt also
\[
E\Bigl(\frac{nt_1}{q}\Bigr) + E\Bigl(\frac{(q-n)t_1}{q}\Bigr) = t_1 - 1,
\]
und daraus durch Multiplication mit $t_1'$ und Summation über alle Werthe von $t_1$ von $1$ bis $q-1$
\[
\sum^{t_1} t_1'\,E\Bigl(\frac{nt_1}{q}\Bigr) + \sum^{t_1} t_1'\,E\Bigl(\frac{(q-n)t_1}{q}\Bigr) \equiv -1 \pmod{q}.
\]

Folglich können sich nicht beide Summen auf der linken Seite durch $q$ theilbar sein. Dann zeigt der Ausdruck (15), dass von den beiden Summen $S_n$ und $S_{q-n}$ gewiss eine durch $q$ untheilbar ist.

Dies aber war zu beweisen.

\bigskip\noindent\textbf{\S. 183. Beweis des zweiten Hülfsatzes für ein ungerades $m$}\\[4pt]

Das zweite Lemma, was nach \S. 181 noch zu beweisen ist, ist folgendes:

\begin{satz}
2. Ist $\varepsilon_n$ ein System conjugirter numerischer Einheiten des Körpers $\Omega_m$, das der Bedingung genügt, dass
\begin{equation}\label{eq:183-1}
\varepsilon_n\,\varepsilon_1^{-n} = \mathcal{E}^m
\end{equation}
die $m^{\text{te}}$ Potenz einer Einheit in $\Omega_m$ ist, so sind die $\varepsilon_n$ selbst $m^{\text{te}}$ Potenzen von Einheiten in $\Omega_m$, multiplicirt mit einer Einheitswurzel der Form $r^{nk}$.
\end{satz}

Auch dieser Beweis ist für den Fall eines ungeraden $m$ ganz elementar. Wir betrachten daher hier zunächst diesen Fall.

Nach \S. 178, 1. können wir jede Einheit des Körpers $\Omega_m$ in der Form darstellen:
\begin{equation}\label{eq:183-2}
\varepsilon = \pm r^k\,\mathcal{E}(r),\quad \varepsilon_n = \pm r^{nk}\,\mathcal{E}(r^n),
\end{equation}
worin $\mathcal{E}(r)$ eine reelle Einheit des Körpers $\Omega_m$ ist. Es genügt also $\mathcal{E}(r)$ der Bedingung
\begin{equation}\label{eq:183-3}
\mathcal{E}(r) = \mathcal{E}(r^{-1}),
\end{equation}
woraus sich nach (2) ergiebt:
\begin{equation}\label{eq:183-4}
\varepsilon_1\,\varepsilon_{-1} = \mathcal{E}(r)^2;
\end{equation}
aus (1) aber findet sich, wenn man $n = -1$ setzt,
\[
\varepsilon_1\,\varepsilon_{-1} = \mathcal{E}^m,
\]
d.\,h.\ die linke Seite von (4) ist die $m^{\text{te}}$ Potenz einer Einheit in $\Omega_m$, und es ist also auch
\begin{equation}\label{eq:183-5}
\mathcal{E}(r)^2 = \mathcal{E}^m
\end{equation}
die $m^{\text{te}}$ Potenz einer solchen Einheit.

Da nun $m$ ungerade ist, so lässt sich die ganze rationale Zahl $x$ so bestimmen, dass
\begin{equation}\label{eq:183-6}
2x + m = 1
\end{equation}
wird, und daraus folgt
\[
\mathcal{E}(r) = \mathcal{E}(r)^{2x}\,\mathcal{E}(r)^m.
\]

Wenn wir also nach (5)
\[
\mathcal{E}(r)^x = e(r)
\]
setzen, so folgt
\begin{equation}\label{eq:183-7}
\mathcal{E}(r) = e(r)^m.
\end{equation}

Durch (7) ist aber der Satz 2. bewiesen und damit zugleich für ein ungerades $m$ das ganze Theorem I.

Auch dieser Schluss versagt für ein gerades $m$, weil dann die Gleichung (6) nicht mehr lösbar ist.

\bigskip\noindent\textbf{\S. 184. Vorläufiges über den Fall eines geraden $m$}\\[4pt]

Für den Fall, dass $m$ eine Potenz von $2$ ist, schlagen wir einen ganz anderen Weg ein, über den hier zunächst einige vorläufige Bemerkungen Platz finden mögen.

Wir haben schon im \S. 182 gesehen, dass das erste Lemma bewiesen ist, wenn wir nachweisen können, dass irgend eine Potenz von $\varphi$, deren Exponent nicht durch $q$ theilbar ist, zur Hauptclasse gehört. Dabei ist von den beiden Eigenschaften des Functionals $\varphi$ nur die erste benutzt, dass die $m^{\text{te}}$ Potenz von $\varphi$ in der Hauptclasse enthalten ist.

Nun kennen wir nach den allgemeinen Sätzen in \S. 154 in jedem Körper einen Exponenten $h$, nämlich die Anzahl der Idealclassen, für den $\varphi^h$ zur Hauptclasse gehört, wenn $\varphi$ ein beliebiges Functional in diesem Körper ist.

Wenn wir also beweisen können, dass die Classenzahl $h$ im Körper $\Omega_m$, wenn $m$ eine Potenz von $2$ ist, eine ungerade Zahl ist, so ist damit ohne alles Weitere das Lemma 1. bewiesen.

Dies soll das Ziel unserer Betrachtungen in den nächsten Abschnitten sein.

In Bezug auf das zweite Lemma liegt die Sache ähnlich. Hier kann man aus den Voraussetzungen in \S. 183, 2. ganz wie oben die Formel \S. 183, (5) herleiten, aus der aber nur zu schliessen ist, dass $\mathcal{E}(r)$ die $\frac{1}{2}m^{\text{te}}$ Potenz einer reellen Einheit $e(r)$ ist.

Nehmen wir also das erste Lemma als bewiesen an, so ergiebt die Formel \S. 181, (14):
\begin{equation}\label{eq:184-1}
(r^n, x_0) = \varrho_n\,\sqrt{e(r^n)}\,\alpha_n \prod (r^{m_2 n}, \eta),
\end{equation}
worin $\varrho_n$ eine Einheitswurzel vom Grade $m^2$, und $\alpha_n$ eine Zahl in $\Omega_m$ ist, und es wäre noch zu beweisen, dass $e(r^n)$ das Quadrat einer Einheit in $\Omega_m$ ist.

Da $e(r)$ eine reelle Einheit ist, so ist
\begin{equation}\label{eq:184-2}
e(r^n) = e(r^{-n}),
\end{equation}
und wir wollen noch festsetzen, dass das Vorzeichen der Wurzel so bestimmt sei (was wir bei passender Annahme über $\varrho_n$ immer annehmen können), dass
\begin{equation}\label{eq:184-3}
\sqrt{e(r^n)} = \sqrt{e(r^{-n})}.
\end{equation}

Das Product $(r^n, x_0)\,(r^{-n}, x_0)$ ist nach der Voraussetzung eine Zahl des Körpers $\Omega_m$, die sich durch die Substitution $(r, r^{-1})$ nicht ändert, d.\,h.\ eine reelle Zahl. Ferner ist nach \S. 174, (5)
\[
(r^{m_2 n}, \eta)\,(r^{-m_2 n}, \eta) = \pm\,p,
\]
also gleichfalls reell. Ebenso ist $\alpha_n\,\alpha_{-n}$ reell, und daraus ergiebt sich nach (1), dass
\[
\varrho_n\,\varrho_{-n} = \frac{(r^n, x_0)\,(r^{-n}, x_0)}{e(r^n)\,\alpha_n\alpha_{-n}\,\prod(r^{m_2 n}, \eta)\,(r^{-m_2 n}, \eta)}
\]
eine reelle Zahl ist, die, weil sie eine Einheitswurzel ist, nur $= \pm 1$ sein kann.

Nun können wir in der hieraus für $n = 1$ sich ergebenden Gleichung
\[
(r, x_0)\,(r^{-1}, x_0) = \pm\,e(r)\,\alpha_1\alpha_{-1}\,\prod(\pm p)
\]
jede Substitution $(r, r^n)$ ausführen, woraus hervorgeht, dass das Zeichen von $\varrho_n\,\varrho_{-n}$ von $n$ unabhängig ist, dass also
\begin{equation}\label{eq:184-4}
\varrho_n\,\varrho_{-n} = \varrho_1\,\varrho_{-1} = \pm 1
\end{equation}
ist. Nun leiten wir aus (1) weiter her:
\begin{equation}\label{eq:184-5}
\varrho_n\,\varrho_1^{-1}\,\sqrt{e(r^n)}\,\sqrt{e(r)}^{-n} = \frac{(r^n, x_0)\,(r, x_0)^{-n}}{\alpha_n\alpha_1^{-1}\,\prod(r^{m_2 n}, \eta)\,(r^{m_2}, \eta)^{-n}},
\end{equation}
und die rechte Seite ist eine Zahl des Körpers $\Omega_m$. Die linke Seite zeigt aber, dass es eine Einheit ist, und wir können demnach, wenn $\mathcal{E}(r)$ eine reelle Einheit bedeutet,
\begin{equation}\label{eq:184-6}
\varrho_n\,\varrho_1^{-1}\,\sqrt{e(r^n)}\,\sqrt{e(r)}^{-n} = r^k\,\mathcal{E}(r)
\end{equation}
setzen. Substituiren wir diesen Werth in die Formel (5) und machen die Substitution $(r, r^{-1})$, so ergiebt sich
\begin{equation}\label{eq:184-7}
\varrho_{-n}\,\varrho_{-1}^{-1}\,\sqrt{e(r^n)}\,\sqrt{e(r)}^{-n} = r^{-k}\,\mathcal{E}(r),
\end{equation}
und durch Multiplication von (6) und (7) mit Rücksicht auf (4)
\begin{equation}\label{eq:184-8}
e(r^n)\,e(r)^{-n} = \mathcal{E}(r)^2.
\end{equation}

In unseren Betrachtungen über die Classenzahl wird sich noch der Satz ergeben:

\begin{quote}
Eine Einheit $e(r)$ des reellen Körpers $H_m$, die mit allen ihren Conjugirten positiv ist, ist das Quadrat einer Einheit im Körper $H_m$.
\end{quote}

Die Formel (8) zeigt aber, dass der Einheit $\pm e(r)$ diese Eigenschaft zukommt, und dass daher $e(r)$ (da auch $-1 = i^2$ das Quadrat einer Einheit ist) das Quadrat einer Einheit des Körpers $\Omega_m$ ist.

Damit ist die Frage auf den Beweis der beiden erwähnten Sätze zurückgeführt, der sich als Schlussstein einer langen, aber an interessanten Beziehungen äusserst reichen Kette von Betrachtungen ergiebt, denen die nächsten Abschnitte gewidmet sein sollen.

\bigskip
\begin{center}
\textbf{Dreiundzwanzigster Abschnitt.}\\[4pt]
\textbf{Classenzahl.}
\end{center}
\bigskip

\noindent\textbf{\S. 185. Der Dirichlet'sche Satz über die Einheiten}\\[4pt]

Die Untersuchungen über die Anzahl der Idealclassen in einem algebraischen Körper, die uns nun die nothwendigen Ergänzungen der Beweise des vorigen Abschnittes geben sollen, bedienen sich der Methoden, die Dirichlet in die Zahlentheorie eingeführt hat\footnote{Dirichlet, Recherches sur diverses applications de l'analyse infinitésimale à la théorie des nombres. Crelle's Journal, Bd. 19, 21. Dirichlet's Werke, Bd. I (1839, 1840). (Hierher gehören übrigens auch die nachgelassenen Abhandlungen von Gauss, „De nexu inter multitudinem classium etc.`` Gauss' Werke, Bd. II.) Die Anwendung dieser Principien auf die Kreistheilungszahlen hat zuerst Kummer gemacht [Crelle's Journal, Bd. 35, 40 (1847, 1850); Liouville's Journal, Bd. 16 (1851)]. Die allgemeine Theorie der Einheiten ist gleichfalls von Dirichlet begründet [Dirichlet's Werke, Bd. I, S. 622, 633, 639 (1840, 1842, 1846)]. Die allgemeine Theorie der Einheiten und die Bestimmung der Classenzahlen hat Dedekind gegeben: Dirichlet-Dedekind, Vorlesungen über Zahlentheorie, 4. Auflage, \S. 183 f. Minkowski, Geometrie der Zahlen.}.

Sie sind nicht mehr rein algebraisch, insofern dabei auch transcendente Functionen und Zahlen, wie der natürliche Logarithmus und die Zahl $\pi$ vorkommen. Auch wird von Grenzübergängen, wie in der Integralrechnung, Gebrauch gemacht, aus denen wir schon in \S. 156 ein schönes Resultat gezogen haben.

Die Untersuchungen, von denen wir zunächst zu handeln haben, sind ganz allgemein, d.\,h.\ sie gelten für jeden algebraischen Zahlkörper, und es bietet keinen Vortheil der Einfachheit, sie auf specielle Körper, wie etwa die Kreistheilungskörper, zu beschränken.

\bigskip\noindent\textbf{Die Einheiten}\\[4pt]

Es sei also $\Omega$ ein Körper $n^{\text{ten}}$ Grades, und
\begin{equation}\label{eq:185-1}
\Omega_1,\; \Omega_2,\; \ldots,\; \Omega_n
\end{equation}
die conjugirten Körper. Die irreducible Gleichung $n^{\text{ten}}$ Grades, deren Wurzeln uns diese conjugirten Körper bestimmen, wird im Allgemeinen sowohl reelle als imaginäre Wurzeln haben, von denen aber die letzteren immer paarweise vorkommen, so dass zu jedem imaginären Körper $\Omega^{(i)}$ ein anderer gehört, dessen Zahlen mit denen von $\Omega^{(i)}$ conjugirt imaginär sind, der also aus $\Omega^{(i)}$ durch die Substitution $(i, -i)$ hervorgeht.

Solche imaginäre Paare fassen wir zu einer Einheit zusammen und bezeichnen die Anzahl der reellen Körper und der Paare imaginärer Körper der Reihe (1) mit $\nu$. Sind alle conjugirten Körper reell, so ist $n = \nu$; sind sie alle imaginär, so ist $n = 2\nu$.

Ausser im Falle des rationalen Körpers ist $\nu$ nur noch in dem Falle eines imaginären quadratischen Körpers $= 1$. Sonst ist $\nu$ immer grösser als $1$.

Die Anzahl der imaginären Paare ist $n - \nu$, und die Anzahl der reellen Körper $2\nu - n$.

In der Determinante, durch die in \S. 145 die Quadratwurzel aus der Grundzahl, $\sqrt{\Delta}$, definirt ist, vertauschen sich durch die Substitution $(i, -i)$ je zwei conjugirt imaginäre Reihen, und $\sqrt{\Delta}$ wechselt also sein Zeichen dann, wenn diese Zahl ungerade ist, sonst nicht. Im ersten Falle ist $\sqrt{\Delta}$ imaginär, im zweiten reell. Daraus ergiebt sich also, dass die Grundzahl $\Delta$ des Körpers $\Omega$ positiv oder negativ ist, je nachdem $n - \nu$ gerade oder ungerade ist.

Behalten wir von zwei conjugirt imaginären Körpern nur den einen bei, so ergiebt sich die Reihe der conjugirten Körper
\[
\Omega_1,\; \Omega_2,\; \ldots,\; \Omega_\nu,
\]
denen wir ein Zeichensystem
\[
\delta_1,\; \delta_2,\; \ldots,\; \delta_\nu
\]
mit der Bestimmung zuordnen, dass $\delta_s = 1$ sein soll, wenn $\Omega_s$ reell, und $= 2$, wenn $\Omega_s$ imaginär ist, also $\sum\delta_s = n$.

Es möge $\eta$ eine von Null verschiedene Zahl des Körpers $\Omega$ bedeuten und
\[
\eta_1,\; \eta_2,\; \ldots,\; \eta_n
\]
die conjugirten Zahlen. Diesem Zahlensysteme ordnen wir ein anderes Zahlensystem zu
\[
\lambda_1,\; \lambda_2,\; \ldots,\; \lambda_\nu,
\]
das wir durch die Gleichung
\begin{equation}\label{eq:185-2}
\lambda_s = \delta_s \log|\eta_s|
\end{equation}
definiren, worin $|\eta_s|$ den absoluten Werth von $\eta_s$ und $\log|\eta_s|$ den reellen natürlichen Logarithmus der positiven Zahl $|\eta_s|$ bedeutet. Ist $\Omega_s$ reell und das Zeichen $\pm$ so gewählt, dass $\pm\eta_s$ positiv ist, so ist
\[
\lambda_s = \log(\pm\eta_s).
\]

Ist aber $\eta_s, \eta_s'$ ein imaginäres Paar, so ist
\[
\lambda_s = \log \eta_s\eta_s',
\]
und zu $\eta_s$ und $\eta_s'$ gehört dasselbe $\lambda_s$, so dass die Anzahl der verschiedenen $\lambda_s$ gleich $\nu$ ist. Aus dieser Bestimmung ergiebt sich noch
\begin{equation}\label{eq:185-3}
\lambda_1 + \lambda_2 + \cdots + \lambda_\nu = \log N_a(\eta),
\end{equation}
wenn, wie früher, $N_a$ die absolute Norm bedeutet.

Die Zahlen $\lambda_1, \lambda_2, \ldots, \lambda_\nu$ wollen wir die \emph{conjugirten Logarithmen} der Zahl $\eta$ nennen.

Wenn $\eta$ eine ganze Zahl des Körpers $\Omega$ ist, so ist die absolute Norm eine natürliche Zahl, die nur dann $= 1$ ist, wenn $\eta$ eine Einheit ist, und daraus folgt nach (3):

\begin{satz}
1. Die Summe der conjugirten Logarithmen einer ganzen Zahl $\eta$ ist positiv, und nur dann gleich Null, wenn $\eta$ eine Einheit ist.
\end{satz}

Bedeutet $\omega_1, \omega_2, \ldots, \omega_n$ eine Basis der ganzen Zahlen in $\Omega$ (eine Minimalbasis von $\Omega$, \S. 145), und $\omega_{1,s}, \omega_{2,s}, \ldots, \omega_{n,s}$ für $s = 1, 2, \ldots, n$ die conjugirten Zahlen, so lassen sich alle ganzen Zahlen des Körpers $\Omega_s$ in der Form darstellen:
\begin{equation}\label{eq:185-4}
\eta_s = x_1\omega_{1,s} + x_2\omega_{2,s} + \cdots + x_n\omega_{n,s},
\end{equation}
mit ganzen rationalen Coëfficienten $x_1, x_2, \ldots, x_n$. Die Determinante dieses Systemes
\[
\sum \pm \omega_{1,1}\,\omega_{2,2} \ldots \omega_{n,n} = \sqrt{\Delta}
\]
ist die Quadratwurzel aus der Discriminante $\Delta$ des Körpers, und demnach immer von Null verschieden. Demnach lässt sich das System (4) nach den Unbekannten $x_1, x_2, \ldots, x_n$ auflösen, und wenn wir nun festsetzen, dass die absoluten Werthe der Zahlen $\eta_s$ nicht über eine gewisse Grenze hinausgehen sollen, so ergeben sich aus diesen Auflösungen Grenzen für die ganzen rationalen Zahlen $x_i$.

Diese einfache Bemerkung liefert uns den folgenden wichtigen Satz:

\begin{satz}
2. Es giebt in einem algebraischen Körper $\Omega$ nur eine endliche Anzahl ganzer Zahlen $\eta$ von der Beschaffenheit, dass die absoluten Werthe der conjugirten Zahlen $\eta_s$ unter einer gegebenen Grenze liegen.
\end{satz}

Wir machen jetzt von dem allgemeinen Satze \S. 155, (25) Gebrauch, nach dem, wenn $f(x_1, x_2, \ldots, x_n)$ irgend eine positive quadratische Form von $n$ Variablen mit der Determinante $D$ ist, die Variablen $x_i$ sich als ganze rationale Zahlen, nicht alle verschwindend, so bestimmen lassen, dass
\[
f(x_1, x_2, \ldots, x_n) < n\sqrt[n]{0{,}404\,D},
\]
d.\,h.\ kleiner als eine von der Determinante $D$ allein abhängige Zahl wird.

Diesen Satz wenden wir, wie im \S. 156, auf die Form an:
\[
f(x_1, x_2, \ldots, x_n) = \frac{|\eta_1|^2}{c_1^2} + \frac{|\eta_2|^2}{c_2^2} + \cdots + \frac{|\eta_n|^2}{c_n^2},
\]
worin die $\eta_s$ die Linearformen (4) sind, und im Falle eines reellen $\eta_s$
\[
|\eta_s|^2 = (\eta_s)^2,
\]
im Falle eines imaginären Paares $\eta_s, \eta_s'$
\[
|\eta_s|^2 = |\eta_s'|^2 = \eta_s\eta_s'.
\]

Im letzteren Falle ist $\eta_s\eta_s'$ die Summe zweier reeller Quadrate, und demnach ist, wenn die $c_s$ reell angenommen werden, $f$ eine positive Form. Wir wollen über die bis jetzt willkürlichen reellen Grössen $c_s$ noch festsetzen, dass, wenn $\eta_s, \eta_s'$ ein conjugirt imaginäres Paar bilden,
\begin{equation}\label{eq:185-5}
c_s = c_s'
\end{equation}
sein soll; ausserdem wollen wir die $c_s$ noch der Bedingung unterwerfen:
\begin{equation}\label{eq:185-6}
c_1 c_2 \ldots c_n = 1.
\end{equation}

Die Determinante der Function $f$ bilden wir, wie im \S. 156, dadurch, dass wir sie zunächst als Form der Variablen $\eta_1, \eta_2, \ldots, \eta_n$ auffassen, und die Determinante dieser Form ist wegen (6) gleich $\pm 1$. Die Determinante von $f$ in Bezug auf die Variablen $x_i$ ist also (Bd. I, \S. 56), vom Vorzeichen abgesehen, gleich dem Quadrate der Determinante der Linearformen $\eta$, d.\,h.\ der Grundzahl des Körpers $\Omega$.

Daraus folgt, dass man eine von den $c_s$ unabhängige, nur durch den Körper $\Omega$ bestimmte Zahl $A$ finden kann, von der Art, dass, was auch die $c_s$ sein mögen, $f(x_1, x_2, \ldots, x_n)$ für ganzzahlige nicht verschwindende $x$ unter $A$ heruntersinkt. Da $f$ die Summe von positiven Summanden ist, so gilt dasselbe von jedem dieser Summanden, und damit ist der folgende Satz bewiesen:

\begin{satz}
3. Ist $c_s$ ein beliebiges, den Bedingungen (5), (6) genügendes System reeller positiver Zahlen, so giebt es eine durch den Körper $\Omega$ allein bestimmte reelle Zahl $A$ von der Art, dass man immer eine ganze Zahl $\eta$ des Körpers $\Omega$ bestimmen kann, die mit ihren conjugirten Zahlen $\eta_s$ den Bedingungen genügt
\begin{equation}\label{eq:185-7}
|\eta_s| < c_s A.
\end{equation}
\end{satz}

Beiläufig folgt hieraus, dass man in jedem algebraischen Körper $\Omega$, ausgenommen dem rationalen und dem imaginären quadratischen, von Null verschiedene ganze Zahlen finden kann, deren absoluter Werth unter jede gegebene Grenze heruntersinkt.

Wir wollen den Ausdruck des Satzes 3. durch Einführung der conjugirten Logarithmen von $\eta$ etwas umformen. Wir setzen $\log A = k$, so dass $k$ eine durch $\Omega$ bestimmte reelle Zahl ist, ferner
\begin{equation}\label{eq:185-8}
c_s = e^{\frac{\gamma_s u}{\delta_s}},
\end{equation}
worin $u$ eine beliebige reelle Grösse sein kann, und die Zahlen $\gamma_s$ wegen (5) und (6) der Bedingung genügen:
\begin{equation}\label{eq:185-9}
\gamma_1 + \gamma_2 + \cdots + \gamma_\nu = \sum_{1,\nu} \gamma_i = 0.
\end{equation}

Da überdies
\[
|\eta_s| = e^{\frac{\lambda_s}{\delta_s}}
\]
ist, so folgt aus (7) und (8):
\begin{equation}\label{eq:185-10}
\lambda_s - \gamma_s u < k\,\delta_s.
\end{equation}

Die Summe der $\nu$ Ausdrücke auf der linken Seite dieser Ungleichung ist wegen (9) gleich $\sum\lambda_s$, also nach dem Satze 1. nicht negativ, und wir finden
\[
0 \lessgtr \sum(\lambda_s - \gamma_s u) < k\,n.
\]

Daraus ergiebt sich wegen (10), dass ein Glied dieser Summe, $\lambda_s - \gamma_s u$, nicht kleiner sein kann, als $-(n - \delta_s)\,k$, weil sonst die Gesammtsumme negativ wäre, und wenn wir der Einfachheit halber die untere Grenze noch etwas kleiner nehmen
\begin{equation}\label{eq:185-11}
-nk\,\delta_s < \lambda_s - \gamma_s u < k\,\delta_s.
\end{equation}

Damit ist bewiesen, dass es für jedes gegebene System der Zahlen $u, \gamma_s$, wenn nur die Bedingung (9) erfüllt ist, eine ganze Zahl $\eta$ des Körpers $\Omega$ giebt, die mit ihren conjugirten Zahlen den Grenzbedingungen (11) genügt.

Es sei nun ferner $g_s$ ein System reeller Grössen, von dem wir nur verlangen, dass
\[
\gamma_1 g_1 + \gamma_2 g_2 + \cdots + \gamma_\nu g_\nu = \alpha
\]
von Null verschieden sei. Damit ist [nach (9)] ausgeschlossen, dass alle $g_s$ einander gleich sind; wenn sie aber das nicht sind, so werden sich zu jedem gegebenen Systeme der $g_s$ unendlich viele Systeme $\gamma_s$ bestimmen lassen, die der Forderung (9) genügen.

Wird nun $k\sum\delta_s g_s = g$ gesetzt, so ergiebt sich aus (11) durch Multiplication mit $g_s$ und Summation
\begin{equation}\label{eq:185-12}
-ng + \alpha u < \sum g_s\lambda_s < g + \alpha u.
\end{equation}

Nach dieser Grenzbedingung bestimmen wir nun, bei festgehaltenen $g_s$ und $\alpha$, eine Reihe von ganzen Zahlen
\begin{equation}\label{eq:185-13}
\eta,\; \eta',\; \eta'',\; \eta''',\; \ldots,
\end{equation}
indem wir für $u$ eine Reihe von Annahmen machen:
\[
u, u', u'', u''', \ldots,
\]
die folgendermaassen näher bestimmt sind:

Wir wählen $u$ zunächst so, dass
\[
-ng + \alpha u = a,\quad g + \alpha u = a'
\]
positiv werden; dann setzen wir
\[
\delta = \frac{a' - a}{\alpha} = \frac{(n+1)\,g}{\alpha},
\]
und setzen
\begin{alignat*}{3}
u' &= u + \delta, &\quad u'' &= u' + \delta, &\quad u''' &= u'' + \delta,\; \ldots,\\
a + \alpha\delta &= a', &\quad a' + \alpha\delta &= a'', &\quad a'' + \alpha\delta &= a''',\; \ldots,
\end{alignat*}
und erhalten so aus (12), wenn $\lambda_s', \lambda_s'', \ldots$ die conjugirten Logarithmen von $\eta', \eta'', \ldots$ sind:
\begin{equation}\label{eq:185-14}
\begin{aligned}
a &< \sum g_s\lambda_s &&< a'\\
a' &< \sum g_s\lambda_s' &&< a''\\
a'' &< \sum g_s\lambda_s'' &&< a'''\\
\cdot\;\cdot\;\cdot\; &\quad\cdot\;\cdot\;\cdot\;\cdot\;\cdot\;\cdot\;\cdot\;\cdot\;\cdot\;\cdot\;\cdot\; &&\cdot\;\cdot\;\cdot\;\cdot\;\cdot\;\cdot\;\cdot\;\cdot\;\cdot
\end{aligned}
\end{equation}
und die Reihe dieser Ungleichungen lässt sich unbegrenzt fortsetzen.

Alle diese Zahlen $\eta, \eta'\eta'', \ldots$ haben überdies nach (6) und (7) die Eigenschaft, dass ihre absoluten Normen $N, N', N'', \ldots$ unter einer bestimmten endlichen Grenze $e^{nk}$ bleiben.

Die Anzahl der möglichen Werthe von $N, N', N'', \ldots$ ist also endlich, nämlich gewiss nicht grösser, als die grösste in $e^{nk}$ enthaltene ganze Zahl. Ebenso ist die Anzahl aller möglichen Reste der Zahlen $x_1, x_2, \ldots, x_n$ [in (4)] nach allen Moduln $N, N', N'', \ldots$ nur eine endliche, und daraus folgt, dass, wenn wir die Reihe der Zahlen (13) nur weit genug fortsetzen, in der Reihe zwei verschiedene Zahlen $\eta^{(h)}, \eta^{(k)}$ auftreten müssen, in denen $N^{(h)} = N^{(k)}$ und die Reste von $x_1, x_2, \ldots, x_n$ nach dem Modul $N^{(h)}$ genau dieselben sind.

Aus (14) aber folgt, wenn $h > k$ vorausgesetzt wird:
\begin{equation}\label{eq:185-15}
\sum g_s\,(\lambda_s^{(h)} - \lambda_s^{(k)}) > 0.
\end{equation}

Ist $N$ die absolute Norm dieser beiden Zahlen $\eta^{(h)}, \eta^{(k)}$, so ist wegen der Uebereinstimmung der Reste der $x$ die Differenz $\eta^{(h)} - \eta^{(k)}$ durch $N$ theilbar. Andererseits ist nach \S. 138 (13) die Zahl $N$ durch $\eta^{(h)}$ theilbar, und folglich ist auch $\eta^{(h)} - \eta^{(k)}$ durch $\eta^{(h)}$ theilbar. Daraus ergiebt sich, dass auch $\eta^{(k)}$ durch $\eta^{(h)}$ und ebenso $\eta^{(h)}$ durch $\eta^{(k)}$ theilbar ist. Beide Zahlen sind also associirt, und wenn wir
\[
\frac{\eta^{(h)}}{\eta^{(k)}} = \varepsilon
\]
setzen, so ist $\varepsilon$ eine Einheit des Körpers $\Omega$.

Setzen wir ausserdem, indem wir mit $|\varepsilon_s|$ den absoluten Werth von $\varepsilon_s$ bezeichnen,
\begin{equation}\label{eq:185-16}
\delta_s \log|\varepsilon_s| = l_s,
\end{equation}
so ergiebt sich aus (15):
\begin{equation}\label{eq:185-17}
g_1 l_1 + g_2 l_2 + \cdots + g_\nu l_\nu > 0.
\end{equation}

\bigskip\noindent\textbf{\S. 186. Systeme unabhängiger Einheiten und Exponentensysteme der Einheiten}\\[4pt]

Dies beweist nun der Satz:

\begin{satz}
4. Ist $g_1, g_2, \ldots, g_\nu$ ein beliebiges System reeller Zahlen, die nicht alle einander gleich sind, so lässt sich in jedem algebraischen Körper $\Omega$ eine Einheit $\varepsilon$ so bestimmen, dass die Summe
\[
g_1 l_1 + g_2 l_2 + \cdots + g_\nu l_\nu
\]
von Null verschieden ist, wenn $l_s$ das System der conjugirten Logarithmen von $\varepsilon$ ist.
\end{satz}

Dieser wichtige Satz, der das Fundament für das Folgende ist, rührt, wenn auch in etwas veränderter Fassung, von Dirichlet her. Diese Formulirung ist von Minkowski gegeben.

Wenn wir in dem zuletzt bewiesenen Satze $g_1 = 1, g_2 = 0, \ldots, g_\nu = 0$ setzen, was, wenn wir von jetzt an $\nu > 1$ voraussetzen, gestattet ist, so folgt, dass es in $\Omega$ eine Einheit $\varepsilon$ giebt, für die einer der conjugirten Logarithmen, $l_1$, von Null verschieden ist. Da jede Potenz von $\varepsilon$ dieselbe Eigenschaft hat, so giebt es auch unendlich viele solche Einheiten.

Dieser Satz gestattet nun eine sehr wichtige Verallgemeinerung:

\begin{satz}
5. Ist $s \lessgtr \nu - 1$, so lässt sich im Körper $\Omega$ ein System von Einheiten mit den zugehörigen conjugirten Logarithmen
\begin{equation}\label{eq:186-1}
\begin{array}{llll}
\varepsilon_1 : & l_{1,1}, & l_{1,2}, & \ldots,\; l_{1,\nu}\\
\varepsilon_2 : & l_{2,1}, & l_{2,2}, & \ldots,\; l_{2,\nu}\\
\multicolumn{4}{c}{\cdot\;\cdot\;\cdot\;\cdot\;\cdot\;\cdot\;\cdot\;\cdot\;\cdot\;\cdot\;\cdot\;}\\
\varepsilon_s : & l_{s,1}, & l_{s,2}, & \ldots,\; l_{s,\nu}
\end{array}
\end{equation}
derart bestimmen, dass eine beliebige der $s$-reihigen Determinanten der Matrix der $l_{h,k}$, etwa
\[
L_s = \sum \pm l_{1,1}\,l_{2,2} \ldots l_{s,s}
\]
von Null verschieden ist.
\end{satz}

Für $s = 1$ ist der ausgesprochene Satz nach der am Eingange gemachten Bemerkung richtig. Wir nehmen also an, er sei bewiesen für $s - 1$ und leiten ihn durch vollständige Induction allgemein her. Dazu ordnen wir die Determinante $L_s$ nach den Elementen der letzten Zeile, und schreiben sie so:
\begin{equation}\label{eq:186-2}
L_s = g_1 l_{s,1} + g_2 l_{s,2} + \cdots + g_\nu l_{s,\nu},
\end{equation}
worin $g_s = L_{s-1}$ ist, und daher nach der gemachten Voraussetzung von Null verschieden angenommen werden kann. Substituiren wir diese Werthe von $g$ in die Formel des Satzes 4, \S. 185, indem wir $g_{s+1} = 0, \ldots, g_\nu = 0$ annehmen, während $g_s$ nicht $= 0$ ist, so sind, so lange $s < \nu$ ist, gewiss nicht alle $g_1, g_2, \ldots, g_\nu$ einander gleich, und es ergiebt sich, dass sich $\varepsilon_s$ so annehmen lässt, dass $L_s$ von Null verschieden ist, wie bewiesen werden sollte.

Auf $s = \nu$ ist diese Schlussweise nicht auszudehnen, und die Determinante $L_\nu$ muss auch in der That immer Null sein, weil für jede Einheit $\sum\limits_{1,\nu} l_s$ verschwindet.

Wir wollen jetzt für den Fall, dass $s = \nu - 1$ ist, die Determinante $L_s$ so bezeichnen:
\begin{equation}\label{eq:186-3}
L_{\nu-1} = L(\varepsilon_1, \varepsilon_2, \ldots, \varepsilon_{\nu-1}),
\end{equation}
und ein System von Einheiten $\varepsilon_1, \varepsilon_2, \ldots, \varepsilon_{\nu-1}$, für welches diese Determinante von Null verschieden ist, ein \emph{System unabhängiger Einheiten} nennen.

Der absolute Werth $L$ der Determinante $L(\varepsilon_1, \varepsilon_2, \ldots, \varepsilon_{\nu-1})$ heisst der \emph{Regulator} des Systems $\varepsilon_1, \varepsilon_2, \ldots, \varepsilon_{\nu-1}$\footnote{Dirichlet-Dedekind, Vorlesungen über Zahlentheorie 4. Auflage, \S. 183.}.

Wenn wir in der Determinante
\begin{equation}\label{eq:186-4}
\begin{vmatrix}
l_{1,1}, & l_{1,2}, & \ldots, & l_{1,\nu}\\
\cdot\;\cdot\;\cdot\;\cdot\;\cdot\;\cdot\;\cdot\;\cdot\;\cdot\;\cdot\;\cdot\;\cdot\;\cdot\;\cdot\;\cdot\\
l_{\nu-1,1}, & l_{\nu-1,2}, & \ldots, & l_{\nu-1,\nu}\\
x_1, & x_2, & \ldots, & x_\nu
\end{vmatrix},
\end{equation}
in der die $x_1, x_2, \ldots, x_\nu$ willkürliche Grössen sind, alle Colonnen zu der letzten addiren, so erhalten wir mit Rücksicht auf die Relationen $\sum\limits_{1,\nu} l_{s,i} = 0$ ihren Werth gleich
\[
\pm(x_1 + x_2 + \cdots + x_\nu)\,L.
\]

Wenn wir also alle $x$ mit Ausnahme von einem beliebigen

\end{document}
